% \documentclass{article}
  % \usepackage{maa-monthly}
 \documentclass{amsart}
% \documentclass{siamltex}


  \let\oldchapter=\chapter
\def\resetpage{\setcounter{page}{1}}
\def\chapter{\expandafter\resetpage\oldchapter}  
 
\usepackage{mathdots}

%\usepackage[all]{xy}

\usepackage{bm}
\usepackage{bigints}
\usepackage{epsfig}
\usepackage{epstopdf}

\usepackage{srcltx,amssymb}

\usepackage{graphicx}

    \usepackage{amsthm}
    \raggedbottom
      \usepackage{amsmath}
    %% IF YOU HAVE FONTS INSTALLED
    %\usepackage{mtpro2}
    \usepackage{amscd}
    %\usepackage{mathtime}\bullet
    \usepackage[title]{appendix}
    \usepackage{slashed}
    \numberwithin{equation}{section} 
       
   %\theoremstyle{theorem}
   \newtheorem{theorem}{Theorem}
    \newtheorem*{thm2}{Proposition A}
    \newtheorem*{thm5}{Proposition A$'$}
    \newtheorem*{thm3}{Proposition B}
    \newtheorem*{thm4}{Proposition B$'$}
  %  \newtheorem{corollary}[theorem]{Corollary}

\newtheorem{remark}[theorem]{Remark}

\newtheorem*{conjecture}{Conjecture}

   \usepackage{appendix}
   
 \newtheorem{lemma}[theorem]{Lemma}
  \newtheorem{corollary}[theorem]{Corollary}
  %  \newtheorem{theorem}{Theorem}
     \newtheorem{proposition}[theorem]{Proposition}
     
    \theoremstyle{definition}
   \newtheorem{definition}[theorem]{Definition}
  
    
    
   \numberwithin{theorem}{section} 
    
  %  \newtheorem{definition}[theorem]{Definition}
  
  \setcounter{tocdepth}{1}
    
    \begin{document}
    
    
    
    \title{Anti-rotors and algebra invariants}
    
       \author{Fred Greensite}
 
       \begin{abstract} We derive what are believed to be three previously unrecognized families of isomorphism invariants for real finite-dimensional unital associative algebras, and explore various implications regarding analysis of algebra structure. 
   \end{abstract}
   
    \maketitle
    
    \markboth{F. Greensite }{Anti-rotors and algebra invariants}

\section{Introduction}
The general goal of classifying algebras with respect to isomorphism equivalence is deemed to be hopeless according to its recognition as a so-called wild problem, e.g., \cite{eom:2024,drozd:1972,belitskii:2005}.  An ironic interpretation of this limitation is that exploration of algebra structure is necessarily an open-ended prospect.  
The usual approaches proceed from consideration of the category of modules over an algebra, such as with algebraic deformation theory as based on Hochschild \mbox{cohomology} and Gerstenhaber algebras \cite{hochschild:1945,gerstenhaber:1964}.  Though there will ultimately be an interface with Hochschild cohomology, we are going to take a very different path, applicable at least to real finite-dimensional unital associative algebras.  



Unless stated otherwise, use of the word ``algebra" always refers to a unital associative algebra with vector space of elements $\mathbb{R}^n$, $n\ge 1$.  The standard topology on $\mathbb{R}^n$ is assumed, as well as the standard basis with elements taken to be column vectors.  Superscript $^T$ indicates dual vector or matrix transpose as applicable, and ``$\,\cdot\,$" is the usual dot product.


\section{The Skewer and the anti-rotor} 
Our program entails examination of the unary operations implied by an algebra's bilinear product operation as a means of studying the algebra itself.
Each such operation, $f:D(A)\rightarrow A$, is a vector field, where $D(A)$ is an appropriate domain in the algebra's vector space of elements.  In examination of a vector field it is of interest to consider the associated infinitesimal rotation tensor (the exterior derivative of the differential 1-form dual to the vector field; its Hodge dual is the curl of the vector field).  Its computation requires choice of a metric (i.e., a symmetric bilinear form, also referred to as a scalar product or metric tensor).  Specification of the standard topology as above entails imposition of the Euclidean inner product as the metric, allowing the computation to proceed.  However, this choice of metric is arbitrary, and can only be considered a default against which the results of all other choices should be compared.  Thus, with respect to a specific choice of unary operation $f$ to study, we are led to consideration of the mapping, \begin{equation}\label{1/23/25.1} \slashed{\text{S}}_f: \mathcal{M} \rightarrow \delta\text{Rot},\end{equation}  
with, \begin{equation}\label{1/23/25.2} \slashed{\text{S}}_f[\mu] = d\left([f(s)]^T \mu \mathfrak{d}s\right).  \end{equation}
 In (\ref{1/23/25.1}), $\mathcal{M}$ is the linear space of all possible metric tensors relevant to $\mathbb{R}^n$ (inclusive of metrics that are not necessarily positive definite and not necessarily nondegenerate), and $\delta\text{Rot}$ is the linear space of possible infinitesimal rotation tensors relevant \mbox{to $\mathbb{R}^n$}.  Equation (\ref{1/23/25.2}) is a formal expression interpreted with respect to the standard basis as follows.  Writing an algebra element as $s=(x_1,\dots,x_n)^T$, we define $\mathfrak{d}(s) \equiv (dx_1,\dots,dx_n)^T$. Thus, $\mu$ is a real symmetric matrix acted upon on the left by the row array $[f(s)]^T$ and on the right by column array $\mathfrak{d}s$.  On the right-hand-side of (\ref{1/23/25.2}), $d(\cdot)$ represents the exterior derivative operation.  Hence, via the expected symbol manipulations, the right-hand-side of (\ref{1/23/25.2}) yields the infinitesimal rotation tensor associated with the vector field $f$ with respect to the metric $\mu$ (or, equivalently, the infinitesimal rotation tensor of vector field $\mu f(s)$ with respect to the Euclidean metric).

$\slashed{\text{S}}_f$ is a linear transformation taking symmetric tensors to skew-symmetric tensors. Consequently, we call $\slashed{\text{S}}_f$ the ``Skewer" (the slash through ``symmetric" S represents a skewer, as with a shish-kabob). Since we assume the standard basis, $\mathcal{M}$ can be represented by a parametrized symmetric $(n\times n)$-matrix with $\frac{n(n+1)}{2}$ independent parameters, i.e., it is given as a symmetric matrix with a different real parameter in each of its upper triangular entries.  Every possible metric tensor is given by some realization of the set of parameters (i.e., a realization results from any particular choice of $\frac{n(n+1)}{2}$ real values for the parameters).  Consequently, this parametrized symmetric matrix is mapped by the Skewer to a parametrized skew-symmetric matrix representing the resulting infinitesimal rotation tensors.  Each entry not on the main diagonal of the latter matrix is a rational function of the $n$ components of an element as a member of $\mathbb{R}^n$, and the above $\frac{n(n+1)}{2}$ parameters appear as coefficients of constituent rational function terms which sum to the entry.  

While the image of the Skewer is thereby complicated, the kernel of the Skewer is another matter.  It is simply a subspace of $\mathcal{M}$;  the algebra's vector space of elements has vanished.  Thinking globally in the context of this less cluttered environment, we observe that
{\it the selected operation $f$ produces a mapping from the set of algebras to the set of subspaces of $\mathcal{M}$ - whose analysis, for arbitrary $f$, might have much to say regarding algebra structure.}
This paper is concerned with establishing the validity of that thesis. 
  Since ``kernel of the Skewer" is a somewhat unwieldy term, we will henceforth instead refer to that entity as the ``anti-rotor" - because each of its members annihilates the infinitesimal rotation tensor.
  
Results in Section 4 and Section 8 justify beginning our analysis with the algebra's multiplicative inversion operation, i.e., the particular choice $f(s)=s^{-1}$.  

     \begin{definition}\label{7/21/23.1} With respect to elements $s$ of a given $n$-dimensional algebra $A$ written in the standard basis as $s= (x_1,\dots,x_n)^T$, where these elements are in an open ball in $\mathbb{R}^n$ consisting only of units and whose central point is the multiplicative identity element $\mathbf{1}_A$, an {\it uncurling metric $L$} is a symmetric member of the matrix algebra $M_n(\mathbb{R})$  satisfying, \begin{equation}\label{4/7/23.1} d\left([Ls^{-1}]\cdot \mathfrak{d}s\right) = d\left((s^{-1})^T L \mathfrak{d}s\right) = 0, \end{equation} where $\mathfrak{d}s \equiv (dx_1,\dots,dx_n)^T$.  The real vector space of uncurling metrics is the algebra's {\it anti-rotor}, $\mathfrak{u}_A$.  It may also be denoted by $\mathfrak{u}$ if no confusion arises.    \end{definition}
     
   
     A requisite open ball in the above definition always exists, as can be easily demonstrated by considering the left regular representation of the algebra. 

The following simple result indicates the elegant manner in which anti-rotors transform.

\begin{theorem} \label{6/6/23.1} If $K:A_1\rightarrow A_2$ is an algebra isomorphism, then \begin{equation}\label{5/23/23.1} \mathfrak{u}_{A_1} = K^T \mathfrak{u}_{A_2} K. \end{equation} \end{theorem}

\begin{proof} By hypothesis, a nonsingular linear change of variable expression for the elements of $A_2$ in terms of the elements of $A_1$ can be written as $s_2= Ks_1$. \mbox{For $L_2\in \mathfrak{u}_{A_2}$} we obtain, \begin{equation} 0 = d\left((s_2^{-1})^T L_2 \mathfrak{d}s_2\right) =  d\left(((Ks_1)^{-1})^T L_2 \mathfrak{d}(Ks_1)\right) = d\left((s_1^{-1})^T \big[K^T L_2K\big]\mathfrak{d}s_1\right), \label{9/21/23.1}  \end{equation} where we have used Definition \ref{7/21/23.1}, the commutativity of a differential with a linear transformation, and the isomorphism condition $(Ks_1)^{-1} = K(s_1^{-1})$.   It follows from (\ref{9/21/23.1}) that every uncurling metric  of algebra $A_2$ is associated with an uncurling metric of $A_1$, i.e., from the last equality in (\ref{9/21/23.1}) we have  $K^TL_2K\in\mathfrak{u}_{A_1}$.  An analogous argument employing $K^{-1}$ indicates that for $L_1\in\mathfrak{u}_{A_1}$ we have the associated $(K^{-1})^T L_1 K^{-1} \in \mathfrak{u}_{A_2}$.    Furthermore, $K^T L_2K\in\mathfrak{u}_{A_1}$ is such that $(K^{-1})^T [K^T L_2 K]K^{-1} = (K^T)^{-1} [K^T L_2 K]K^{-1} = L_2$.  Similarly, $(K^{-1})^T L_1K^{-1}\in\mathfrak{u}_{A_2}$ is such that $K^T[(K^{-1})^T L_1K^{-1}] K = L_1$.  The theorem follows from the implied bijection. 
\end{proof}

Computation of $\mathfrak{u}$ is straightforward in principle.  The algorithm is contained in the proof of the following.  Since it will cause no confusion, in the sequel we simply write an algebra element in the standard basis as, e.g., $s=(x_1,\dots,x_n)$, where the latter is understood to be a column vector.

\begin{proposition}\label{9/2/24.1} The members of $\mathfrak{u}$ are the realizations of a symmetric $(n\times n)$-matrix $M_{\mathfrak{u}}$ whose entries are linear combinations of the components of an $m$-tuple of real parameters, where $m$ is the dimension of $\mathfrak{u}$ as a vector space, and where each realization of $M_{\mathfrak{u}}$ as a real matrix results from replacing its $m$ parameters with the members of any set of $m$ members of $\mathbb{R}$.   
\end{proposition}

\begin{proof} For an $n$-dimensional algebra with elements $s = (x_1,x_2,\dots,x_n)$ with respect to the standard basis, one initially takes $\mathfrak{u}$ to be a symmetric $(n\times n)$-matrix $M$ whose upper triangular entries are each a different member of a set of $\frac{n(n+1)}{2}$ real parameters.  Computation of $s^{-1}$ produces a vector whose components are each a rational function of the $x_1,x_2,\dots,x_n$.  So far, $Ms^{-1}$ can be understood as a parametrized \mbox{$n$-dimensional} vector field whose components (which can easily be computed) are dependent on the initial set of parameters in $M$ and the $x_1,x_2,\dots,x_n$.  There are then $\frac{n(n-1)}{2}$ constraint equations that must be satisfied in order that (\ref{4/7/23.1}) hold (for $L=M$).  Each of these equations has the form, \begin{equation}\label{9/2/24.2} \frac{\partial}{\partial x_i}([Ms^{-1}]_j) - \frac{\partial}{\partial x_j}([Ms^{-1}]_i) = 0,\end{equation} where $[Ms^{-1}]_k$ indicates the $k$-th component of $Ms^{-1}$.  As noted, we can write $s^{-1} = (\hat{x}_1,\dots,\hat{x}_n)$, where each $\hat{x}_k$ is a rational function of the components of $s$.  Consequently, $[Ms^{-1}]_k$ is a linear combination of rational functions of the components $\hat{x}_1,\dots,\hat{x}_n$, where the coefficients defining this linear combination are linear combinations of various of the above parameters.  It follows that $\frac{\partial}{\partial x_p}([Ms^{-1}]_q)$ is again a sum rational functions of the components $x_1,\dots,x_n$, where this expression is again linear in the parameters.  Thus, multiplying both sides of (\ref{9/2/24.2}) by a polynomial determined by the products of the polynomial denominators of the latter rational functions, the left-hand-side of (\ref{9/2/24.2}) becomes a polynomial in the components of $s$, while remaining an expression that is linear in the parameters - and the right-hand-side of (\ref{9/2/24.2}) remains zero.
Hence, satisfaction of each of the $\frac{n(n-1)}{2}$ constraint equations (\ref{9/2/24.2}) ultimately reduces to the requirement that a particular multivariate polynomial be {\it identically} zero.  After simplification of the polynomial to be a linear combination of independent monomials, each monomial term is then comprised of a coefficient consisting of a linear combination of some of the above $\frac{n(n+1)}{2}$ parameters, multiplying an expression consisting of a product of some of the components of $s$, each raised to some power.  Therefore, each of the above coefficients of these (independent) monomials are required to be zero.  Accordingly, the number of surviving independent parameters (i.e., parameters not required to be equated to some linear combination of other parameters, or set to 0) is in general reduced as the $\frac{n(n-1)}{2}$ constraint equations are successively treated.  $M$ is then modified as necessary, as each successive constraint equation is treated.  After treatment of all the constraint equations via the above steps defined by treatment of the constraint equations in some sequence, the final number $m$ of surviving independent parameters is the dimension of $\mathfrak{u}$, and the final entries of $M$ evidently yield the requisite $M_\mathfrak{u}$.
\end{proof}

Consequently, $\mathfrak{u}$ can be represented as a symmetric $(n\times n)$-matrix $M_{\mathfrak{u}}$ whose entries are particular linear combinations of $m$ independent parameters (where $m$ is the dimension of $\mathfrak{u}$), because each uncurling metric results from a single set of selected real values for those parameters.  The determinant of $M_{\mathfrak{u}}$, written as ${\rm det}(M_{\mathfrak{u}})$, is a polynomial in some or all of the components of an ordered $m$-tuple of parameters, $(\alpha_1,\dots,\alpha_m)$.  If $m$ is small  we will denote the parameters using different Greek letters like $\alpha$, $\beta$, $\gamma$, etc. 

   With the above in mind, it is seen that (\ref{5/23/23.1}) has immediate consequences - our first family of algebra isomorphism invariants resulting from the anti-rotor. 
           
  \begin{corollary}\label{8/18/23.2} The following are algebra isomorphism invariants. \begin{itemize} \item The dimension $m$ of $\mathfrak{u}$.   \item The largest rank  of a member of $\mathfrak{u}$.  \item The smallest rank of a nonzero member of $\mathfrak{u}$. \item The number of parameters on which ${\rm det}(M_{\mathfrak{u}})$ depends. \item Invariants of the real algebraic set, \[V_\mathbb{R}({\rm det}(M_{\mathfrak{u}})) = \{(\alpha_1,\dots,\alpha_m)\in\mathbb{R}^m: {\rm det}(M_{\mathfrak{u}}) = 0 \}.\] \end{itemize} \end{corollary}


\section{The space dual to the anti-rotor}
\enlargethispage{1cm}
  \begin{definition}\label{9/21/23.2} Let $L$ be an uncurling metric for some algebra $A$.  For $s$ in an open ball centered at the multiplicative identity $\mathbf{1}$ and devoid of non-units, with $\|{\bf 1}\|^2 \equiv \mathbf{1}\cdot\mathbf{1}$, the associated {\it unital norm} $ \ell(s)$ is the function 
  \begin{equation}\label{11/7/20.2}   \ell(s) \equiv \exp\left(\frac{1}{\|{\bf 1}\|^2}\int_{\bf 1}^s [Lt^{-1}]\cdot \mathfrak{d}t\right).\end{equation} 
   The set of unital norms associated with an algebra, which is an abelian group $\mathfrak{u}^d_A$ with respect to the usual point-wise operation of multiplication of functions, is called the {\it unital norm group}, also known as the {\it dual of the anti-rotor}. $\mathfrak{u}^d_A$ may be denoted as $\mathfrak{u}^d$ if no confusion arises. \end{definition} 
   
   
The integral in (\ref{11/7/20.2}) is path-independent due to (\ref{4/7/23.1}).  The duality referred to in Definition \ref{9/21/23.2} is easily appreciated.  Equation (\ref{9/21/23.2}) can be rewritten as  \mbox{$\log\ell(s) = \frac{1}{\|\mathbf{1}\|^2}\int_{{\bf 1}}^{s} \left[L t^{-1}\right]\cdot \mathfrak{d}t$.}   We can view this as a forward integral transform, mapping $L$ to $\log \ell(s)$.   The latter equation implies that an uncurling metric and its associated unital norm satisfy,  \begin{equation}\label{11/7/20.1} Ls^{-1}  = \nabla\left[\|{\bf 1}\|^2\log  \ell(s)\right] =   \left(\frac{\|{\bf 1}\|^2}{ \ell(s)}\right)\nabla \ell(s).\end{equation}  Due to invertibility of the multiplicative inversion operation,  (\ref{11/7/20.1}) can be rewritten as  $Ls = \|\mathbf{1}\|^2 \nabla\left[\log \ell(s^{-1})\right]$, where the right-hand-side is interpreted to mean the product of $\|\mathbf{1}\|^2$ with the gradient of $\log \ell$ which is then evaluated at $s^{-1}$.  The latter equation defines the inverse transform - where $L$ (which can be extracted from $Ls$) is obtained from $\log\ell(s)$ (this being an integral transform in the sense of distributions).  
In other words, groups of uncurling metrics and groups of unital norms are dual (where the duality between individual members of  $\mathfrak{u}_A$ and $\mathfrak{u}_A^d$ follows from an integral transform defined with respect to the multiplicative inversion operation $s\mapsto s^{-1}$).
   
    
Also note that according to (\ref{11/7/20.1}), the pair $(L, \ell(s))$ characterizes the algebra's multiplicative inversion operation $s\mapsto s^{-1}$ in the portion of an open ball centered at $\mathbf{1}$ devoid of non-units that is orthogonal to the kernel of $L$.  If $L$ is nonsingular, this pair fully characterizes the inversion operation $s^{-1}$ in an open ball centered at $\mathbf{1}$ devoid of non-units.  For a given algebra, the set of all such pairs (where $L$ may or may not be singular) is an abelian group with respect to the binary operation defined by,
\begin{equation}\label{5/23/23.2} \big(L_1, \ell_1(s)\big)\big(L_2, \ell_2(s)\big) \equiv 
\big(L_1+L_2, \ell_1(s)  \ell_2(s)\big).\end{equation} 
\begin{definition} For a given algebra, the set of pairs $(L, \ell(s))$ satisfying (\ref{4/7/23.1}) and (\ref{11/7/20.2}), in the context of the product operation defined by (\ref{5/23/23.2}), is the algebra's {\it inversion characterization group}. \end{definition}

Any member of the inversion characterization group that is associated with a nonsingular uncurling metric implies the full group, since the member implies the operation $s\mapsto s^{-1}$ via (\ref{11/7/20.1}), from which the full inversion characterization group can be computed.  Obviously, this group is not the direct product $\mathfrak{u}\times\mathfrak{u}^d$. 

Corollary \ref{8/18/23.2} has already identified invariants associated with $\mathfrak{u}$ (whose members supply the first component of the inversion characterization group), but another set of useful invariants becomes apparent on examination of $\mathfrak{u}^d$ (whose members supply the second component of the inversion characterization group).  Based on the mapping from $\mathfrak{u}$ to $\mathfrak{u}^d$ defined by (\ref{11/7/20.2}),  we can write \begin{equation}  \mathfrak{u}^d = \exp\left(\int_{\bf 1}^s [\mathfrak{u}t^{-1}]\cdot \mathfrak{d}t\right). \label{6/7/23.1}\end{equation}  In the above, the notation $\mathfrak{u}t^{-1}$ means $\{Lt^{-1}: L\in \mathfrak{u}\}$.  Accordingly, the right-hand-side of (\ref{6/7/23.1}) is understood to be a collection of exponentiated integrals.   Note that in (\ref{6/7/23.1}) we do not need the factor $\frac{1}{\|{\bf 1}\|^2}$, since $\frac{1}{\|{\bf 1}\|^2}\mathfrak{u}=\mathfrak{u}$.  

The invariants stemming from the $\mathfrak{u}^d$ contribution to the inversion characterization group proceed from the following simple corollary of Theorem \ref{6/6/23.1}, which we will give ``co-equal status" as a theorem because it ultimately addresses the co-important  partner of $\mathfrak{u}$, i.e., $\mathfrak{u}^d$. 

\begin{theorem}\label{6/6/23.2} Let $K:A_1\rightarrow A_2$ be an algebra isomorphism, where for $s_1\in A_1$ we define $s_2\equiv Ks_1$.  Then for $s_1$ in an open ball centered at $\mathbf{1}_{A_1}$ consisting only of units, \begin{equation}\label{9/6/24.2} \int_{{\bf 1}_{A_1}}^{s_1} \left[\mathfrak{u}_{A_1}t_1^{-1}\right]\cdot \mathfrak{d}t_1 
= \int_{{\bf 1}_{A_2}}^{s_2} \left[\mathfrak{u}_{A_2}t_2^{-1}\right]\cdot \mathfrak{d}t_2 . \end{equation} \end{theorem}

\begin{proof} 
Equation (\ref{5/23/23.1}) implies,
\begin{eqnarray*} \int_{{\bf 1}_{A_1}}^{s_1} \left[\mathfrak{u}_{A_1}t_1^{-1}\right]\cdot \mathfrak{d}t_1 &=& \int_{{\bf 1}_{A_1}}^{s_1} \left[ K^T\mathfrak{u}_{A_2}Kt_1^{-1}\right]\cdot \mathfrak{d}t_1 \\ &=& 
\int_{{\bf 1}_{A_2}}^{s_2} \left[ K^T\mathfrak{u}_{A_2}K (K^{-1}t_2)^{-1}\right]\cdot \mathfrak{d}\left(K^{-1}t_2\right) \\ &=& \int_{{\bf 1}_{A_2}}^{s_2} [K^T\mathfrak{u}_{A_2}t_2^{-1}]\cdot \left(K^{-1}\mathfrak{d}t_2\right) = \int_{{\bf 1}_{A_2}}^{s_2} \left[\mathfrak{u}_{A_2}t_2^{-1}\right]\cdot \mathfrak{d}t_2, 
\end{eqnarray*} 
where we have used the isomorphism assumption (which implies $(K^{-1}t_2)^{-1} = K^{-1}t_2^{-1}$), the commutativity of a linear transformation and a differential, and exploitation of an adjoint in the context of an inner product.
\end{proof}

Proposition \ref{9/2/24.1} indicates that $m$-dimensional $\mathfrak{u}$ can be written as a matrix $M_\mathfrak{u}$ whose entries are each a linear combination of the components of an $m$-vector of real parameters $(\alpha_1, \dots, \alpha_m)$.   Thus, $\mathfrak{u}$ can be equivalently thought of as a set of real matrices, each of whose members is given by a choice of real values for the $m$ parameters in the above parametrized matrix version of $\mathfrak{u}$, with each matrix in the set being an uncurling metric.  Now, for each $q\le m$ consider the real matrix  $\mathfrak{u}_q$ obtained from $\mathfrak{u}$ by setting $\alpha_q =1$ with all other parameters set to zero.  These are a set of generators for $\mathfrak{u}$. Thus, $\mathfrak{u}t^{-1} = M_{\mathfrak{u}}(t^{-1}) = \sum_{q=1}^m \alpha_q \mathfrak{u}_q t^{-1}$.   Using (\ref{6/7/23.1}), we then obtain \begin{equation}\label{9/26/23.1} \log  \mathfrak{u}^d = \int_{\mathbf{1}}^s [M_{\mathfrak{u}} (t^{-1})]\cdot \mathfrak{d}t =\sum_{q=1}^m \left(\alpha_q \int_{\mathbf{1}}^s [\mathfrak{u}_q t^{-1}]\cdot \mathfrak{d}t\right).\end{equation}  

We now consider the algebra elements in an open ball centered at $\mathbf{1}$ consisting only of units, with each element written in coordinates as $s = (x_1,\dots,x_n)$, and where $\mathbf{1} = ({\sf 1}_1,{\sf 1}_2,\dots,{\sf 1}_n)$.  The path-independent integral inside the sum on the right-hand-side of (\ref{9/26/23.1}) can be evaluated on the piecewise linear path with successive endpoints of the piecewise linear components as  $({\sf 1}_1,{\sf 1}_2,\dots,{\sf 1}_n)$, $(x_1,{\sf 1}_2,\dots,{\sf 1}_n)$, $(x_1,x_2,{\sf 1}_3,\dots,{\sf 1}_n)$,  $\dots$, $(x_1,\dots, x_{n-1}, {\sf 1}_n)$, $(x_1,x_2,\dots,x_n)$.    On a linear portion of the path, the integrand is a rational function of a single variable.  Thus, on each linear path component the indefinite integral evaluates to a sum of rational functions, logarithms, and arctangents.  However, coefficients in that sum will be functions of the components of $s$ that do not vary on the linear piece of the path.  For dimension less than five, we can obtain a closed form expression for the right-hand-side of (\ref{9/26/23.1}) as a function of the components of $s$.  However, this is not possible in general for dimensions five and greater, given the need for an expression for the roots of a polynomial appearing in the denominator of an integrand - these roots being functions of the nonvarying components of $s$ (of course, if definite values are given for the non-varying components of $s$, the integration on each linear path component can be accomplished to any degree of desired accuracy via well-known polynomial root-finding algorithms).

But while lack of a closed form solution for the roots of a polynomial in dimensions five and above limits the ability to write down a closed form expression for $\log\mathfrak{u}^d$ for an arbitrary $s$ in the given domain, we can evaluate (to any degree of accuracy) the set of $n$ functions given by $\log\mathfrak{u}^d$ computed along paths each beginning at $\mathbf{1}$ and proceeding parallel to each of the component axes of $s$ - because each integrand again has only a single component of $s$ that is varying on the path, but now the other components of $s$ only appear as their known component values in $\mathbf{1}$.  
So, if we successively compute the integrals on the linear path from $\mathbf{1}$ to $(x_1,{\sf 1}_2,\dots,{\sf 1}_n)$, then on the linear path from \mbox{$\mathbf{1}$ to} $({\sf 1}_1,x_2,{\sf 1}_3,\dots,{\sf 1}_n)$, etc., we obtain that set of $n$ functions. 
 Now, evaluation of $q$-th term on the right-hand-side of (\ref{9/26/23.1}) is \begin{equation}\label{7/26/24.1} \alpha_q\int_{\mathbf{1}}^s [\mathfrak{u}_q t^{-1}]\cdot \mathfrak{d}t = \alpha_q\sum_{k=1}^n \int_{\bar{s}_k} [\mathfrak{u}_q t^{-1}]\cdot \mathfrak{d}t,\end{equation} with $\bar{s}_k \equiv [(x_1,\dots,x_{k-1},{\sf 1}_{k},{\sf 1}_{k+1}\dots,{\sf 1}_n),(x_1,\dots,x_{k-1},x_k,{\sf 1}_{k+1},\dots,{\sf 1}_n)]$ being the directed line segment with the indicated endpoints.   As we have noted, evaluation of the integrals inside the sum on the right-hand-side of (\ref{7/26/24.1}) is problematic.  But instead consider,  \begin{equation}\label{7/26/24.2}\alpha_q\sum_{k=1}^n \int_{\hat{s}_k} [\mathfrak{u}_q t^{-1}]\cdot \mathfrak{d}t,\end{equation} with $\hat{s}_k \equiv [({\sf 1}_1,\dots,{\sf 1}_{k-1},{\sf 1}_{k},{\sf 1}_{k+1}\dots,{\sf 1}_n),({\sf 1}_1,\dots,{\sf 1}_{k-1},x_k,{\sf 1}_{k+1},\dots,{\sf 1}_n)]$ being the directed line segment. 
For reasons stated above, approximations to the integrals in (\ref{7/26/24.2}) in terms of rational functions, logarithms, and arctangents, can be obtained with arbitrary accuracy.  Furthermore, for $s$ in a small enough neighborhood of $\mathbf{1}$, the integrands of each integral in (\ref{7/26/24.2}) are close to the integrands appearing in each integral on the right-hand-side of (\ref{7/26/24.1}).  This means that there will be the same number of respective rational function, logarithm, and arctangent terms in the evaluated indefinite integral version of the $k$-th integral in the sum (\ref{7/26/24.2}) as there would be if one evaluated the indefinite integral version of the $k$-th integral in the sum on the right-hand-side of (\ref{7/26/24.1}) (i.e., if any definite values for the non-varying components of $s$ compatible with the neighborhood of $\mathbf{1}$ were inserted into the integrand of (\ref{7/26/24.1})).  Consequently, following evaluation of the integrals, a parameter $\alpha_q$ multiplies a rational function (resp., logarithm, arctangent) in the $k$-th term of the sum in (\ref{7/26/24.2}) if and only if $\alpha_q$ multiplies a rational function (resp., logarithm, arctangent) in the $k$-th term of the sum on the right-hand-side of (\ref{7/26/24.1}) were the integral actually computed for any specific values of the non-varying components of $s$ in the range of values compatible with the small neighborhood of $\mathbf{1}$.  


So, as regards evaluation of the association of parameters of $M_{\mathfrak{u}}$ with rational function, logarithm, and arctangent terms arising from $\log\mathfrak{u}^d$, instead of using the right-hand-side of (\ref{9/26/23.1}) it is sufficient to instead deal with the function,
\begin{equation}\label{7/29/24.1} \sum_{q=1}^m \left(\alpha_q\sum_{k=1}^n \int_{\hat{s}_k} [\mathfrak{u}_q t^{-1}]\cdot \mathfrak{d}t\right).\end{equation} That is, evaluation of (\ref{7/29/24.1}) will reveal the same associations of parameters with rational functions, logarithms, and arctangents, as the right-hand-side of (\ref{9/26/23.1}) in a small neighborhood of $\mathbf{1}$.  Now, following evaluation of the integrals in (\ref{7/29/24.1}), the function expressed by (\ref{7/29/24.1}) is given by respective sums of rational function terms, logarithm terms, and arctangent terms.  Note that for this summed function resulting from algebra $A$, its rational function terms can be collected into a {\it single} rational function term involving some number $\tau_{A,{\rm rat}}$ of distinct parameters.  Furthermore, its logarithm terms can be collected into a {\it single} (real) logarithm involving some number $\tau_{A,{\rm log}}$ of distinct parameters.  Similarly, its arctangent terms, which can be alternatively expressed utilizing $\text{arctan}(x) = \frac{i}{2}\log\frac{1-ix}{1+ix}$, can be collected into a {\it single} such logarithm term utilizing imaginary coefficients, involving some number $\tau_{A,{\rm arc}}$ of distinct parameters. 
Based on dimensional considerations relating to the respective numbers of distinct parameters associated with the single rational function term, the single real logarithm term, and the single logarithm term with imaginary coefficients, it seen that Theorem \ref{6/6/23.2} implies,


\begin{corollary}\label{7/25/24.1} With respect to the set of parameters $\{\alpha_q:q=1,\dots,m\}$ arising in $M_{\mathfrak{u}_A}$, and the function defined by (\ref{7/29/24.1}) following evaluation of the integrals, let $\tau_{A,{\rm rat}}$ be the number of distinct parameters multiplying at least one rational function term, let $\tau_{A,{\rm log}}$ be the number of distinct parameters multiplying at least one logarithm, and let $\tau_{A,{\rm arc}}$ be the number of distinct parameters multiplying at least one arctangent.  If $A_1$ and $A_2$ are isomorphic algebras, then \[ (\tau_{A_1,{\rm rat}},\tau_{A_1,{\rm log}},\tau_{A_1,{\rm arc}}) = (\tau_{A_2,{\rm rat}},\tau_{A_2,{\rm log}},\tau_{A_2,{\rm arc}}).\]
\end{corollary}

The above identifies the second family of algebra invariants resulting from the anti-rotor.


\section{The ``geometries" associated with an algebra} 
On an open ball centered at $\mathbf{1}$ devoid of non-units but containing the line segment from $\left(\frac{1}{\alpha}\right)\mathbf{1}$ to $\alpha\mathbf{1}$ for some choice of $\alpha>0$ (such a neighborhood always exists if $\alpha$ is close enough to 1), it is easy to show from (\ref{11/7/20.2}) that for unital norms in general,  \[  \ell(\alpha s) = \alpha^{\frac{\mathbf{1}^T L \mathbf{1}}{\|\mathbf{1}\|^2}} \ell(s),\] using an argument similar to that employed in the proof of Theorem \ref{9/24/23.1}. Thus, $ \ell(s)$ is a positive homogeneous function of degree given by the exponent on the right-hand-side of the above equation.  
   {\it But things become somewhat nicer if elements that are inverses with respect to the algebra's product also behave inversely with respect to the scalar product defined by $L$}.  
    This situation results if uncurling metric $L$ satisfies the additional constraint (\ref{7/7/21.1}).

    \begin{definition}\label{9/20/23.1} A {\it normalized uncurling metric} $L_{\shortparallel}$ is an uncurling metric that satisfies, \begin{equation} \label{7/7/21.1} s^TL_{\shortparallel} s^{-1} =\|{\bf 1}\|^2,\end{equation} on a sufficiently small neighborhood of $\mathbf{1}$.   A {\it special unital norm} $\ell_{\rm sp}(s)$ is a unital norm that is derived from a normalized uncurling metric $L_{\shortparallel}$ according to (\ref{11/7/20.2}). 
\end{definition}


\begin{theorem}\label{9/24/23.1} $ \ell_{\rm sp}(s)$ is a degree-1 positive homogeneous function, and
\begin{equation}\label{5/23/23.3}  \ell_{\rm sp}(s^{-1}) = \big(\ell_{\rm sp}(s)\big)^{-1} = \frac{1}{\ell_{\rm sp}(s)}, \end{equation} when $s,s^{-1}$ are in an open ball centered at $\mathbf{1}$ devoid of non-units. \end{theorem}

\begin{proof} We have already observed (immediately following Definition \ref{7/21/23.1}) that an open ball centered at $\mathbf{1}$ devoid of non-units always exists.  Let $\eta$  be such a neighborhood.

Given $\frac{1}{\alpha}s,s,\alpha s\in\eta$ for some choice of $\alpha>0$, we have from (\ref{11/7/20.2}) that
\begin{eqnarray}\nonumber  \log\ell_{\rm sp}(\alpha s) &=&  \frac{1}{\|{\bf 1}\|^2}\int_{\bf 1}^{\alpha s} [L_{\shortparallel}t^{-1}]\cdot \mathfrak{d}t = \frac{1}{\|{\bf 1}\|^2}\int_{\frac{\mathbf{1}}{\alpha}\mathbf{1}}^{s} [L_{\shortparallel}(\alpha w)^{-1}]\cdot \mathfrak{d}(\alpha w)  \\ \nonumber &=&  \frac{1}{\|{\bf 1}\|^2}\int_{\mathbf{1}}^{s} [L_{\shortparallel}w^{-1}]\cdot \mathfrak{d}w + \frac{1}{\|{\bf 1}\|^2}\int_{\frac{\mathbf{1}}{\alpha} \mathbf{1}}^{\mathbf{1}} [L_{\shortparallel}w^{-1}]\cdot \mathfrak{d}w  \\ &=& \log \ell_{\rm sp}(s) + \log\alpha, \label{4/2/23.1}
\end{eqnarray} where the integral from $\frac{\mathbf{1}}{\alpha} \mathbf{1}$ to $\mathbf{1}$ can be easily evaluated along the line segment with those endpoints, and we have used (\ref{7/7/21.1}) in evaluation of that integral to  obtain the term $\log \alpha$ on the right-hand-side of the final equation.  Thus, we have $\ell_{\rm sp}(\alpha s) = \alpha\ell_{\rm sp}(s)$, degree-1 positive homogeneity.


Let $\eta^{-1}$ be the set consisting of the multiplicative inverses of all members of $\eta$.  It is easily shown that $\eta^{-1}$ is an open neighborhood of $\mathbf{1}$, and $\eta$, $\eta^{-1}$ are diffeomorphic via the multiplicative inversion operation.  
We now consider a smooth path $\mathcal{P}$ from $\mathbf{1}$ to $s$ which is inside $\eta$.  Let $\mathcal{P}^{-1}$ be the point set consisting of the multiplicative inverses of the members of $\mathcal{P}$.  Clearly, $\mathcal{P}^{-1}$ is contained in $\eta^{-1}$. 
\mbox{We then have,} \begin{eqnarray}  \log \ell_{\rm sp}(s^{-1})   &=&   \frac{1}{\|{\bf 1}\|^2} \int_{\bf 1}^{s^{-1}} [L_{\shortparallel}t^{-1}] \cdot \mathfrak{d}t \nonumber \\ &=&  \frac{1}{\|{\bf 1}\|^2} \int_{\bf 1}^s [L_{\shortparallel}y]\cdot \mathfrak{d}\left(y^{-1}\right) \nonumber \\ &=&  \frac{1}{\|{\bf 1}\|^2}\left(\left(y^{-1}\cdot [L_{\shortparallel}y]\right)\bigg|_{\bf 1}^s - \int_{\bf 1}^s y^{-1}\cdot \mathfrak{d}\left(L_{\shortparallel}y\right)\right)   \nonumber \\ &=& -\frac{1}{\|{\bf 1}\|^2} \int_{\bf 1}^s y^{-1}\cdot (L_{\shortparallel}\mathfrak{d}y) = -\frac{1}{\|{\bf 1}\|^2} \int_{\bf 1}^s [L_{\shortparallel}y^{-1}] \cdot \mathfrak{d}y  \nonumber \\ &=& -\log \ell_{\rm sp}(s),\label{4/2/23.5} \end{eqnarray} where we have used the change of variable $y=t^{-1}$, commutativity of a linear transformation and a differential, the property that $L_{\shortparallel}$ is a real symmetric matrix, and  (\ref{7/7/21.1}).  Equation (\ref{5/23/23.3}) then follows.
\end{proof}

As demonstrated in Section 8 (Theorem \ref{11/17/22.1}), every algebra is associated with at least one special unital norm.

It is also seen in Section 8 that at least some members of an algebra's group of unital norms are multiplicative, i.e., $\ell(s_1s_2) = \ell(s_1)\ell(s_2)$ (in fact, all of the unital norms associated with the two-dimensional algebras are multiplicative).   But in general, unital norms are not multiplicative.  However, the unital norms originate from the nonlinear unary multiplicative inversion operation $s\mapsto s^{-1}$ implied by the algebra's bilinear binary product operation - so the particularly desirable fealty would be to that unary operation.  It is thus significant that the special unital norms pay homage to it via (\ref{5/23/23.3}) (also, see comments in the second and third paragraphs of Section 10.2).

Two reasons we have employed the term {\it unital} in ``unital norm" are that the latter entities apply to {\it unital} algebras, and they {\it unite} metric and rotor information.  One reason we have employed the term ``norm" for these entities is that the group includes a member closely related to the usual algebra norm (Corollary \ref{11/17/22.1}).  A second reason is as follows.   Using (\ref{11/7/20.1}), (\ref{5/23/23.3}), the degree-1 positive homogeneity property of $\ell_{\rm sp}(s)$, and assuming the associated uncurling metric $L_{\shortparallel}$ is nonsingular, it is appreciated that a special unital norm has the feature that, 
\begin{equation}\label{5/22/24.5}  s = \ell_{\rm sp}(s)\Big(\|\mathbf{1}\|^2 L_{\shortparallel}^{-1}[\nabla\ell_{\rm sp}(s^{-1})]\Big),\end{equation} with
\begin{equation}\label{5/22/24.6}\ell_{\rm sp}\Big(\|\mathbf{1}\|^2 L_{\shortparallel}^{-1}[\nabla\ell_{\rm sp}(s^{-1})]\Big) = 1,\,\end{equation}
where $\nabla\ell_{\rm sp}(s^{-1})$ means that $\nabla\ell_{\rm sp}$ is evaluated at $s^{-1}$.  What these equations say is that an element $s$ is the product of its ``norm" $\ell_{\rm sp}(s)$ and an element of unit ``norm" that is determined by the gradient of $\ell_{\rm sp}$ evaluated at a point dependent on $s$ - so that the factor enclosed by large parentheses on the right-hand-side of (\ref{5/22/24.5}) is  ``the direction of $s$".  

Equations (\ref{5/22/24.5}), (\ref{5/22/24.6}) are very reminiscent of equations satisfied by the norm on an inner product space, or more generally, on any quadratic space, via the ``norm" $\ell(s)$ \mbox{specified by} \begin{equation}\label{11/29/23.1} s = \ell(s)L^{-1}[\nabla\ell(s)],\end{equation} and \begin{equation}\label{11/29/23.2}\ell\left(L^{-1}[\nabla\ell(s)]\right) = 1,\end{equation} where $L$ is the scalar product matrix defining the quadratic space. For example, assuming $L$ has at least one positive eigenvalue, in a small neighborhood $\eta$ of a point $\hat{s}$ such that $\hat{s}^T L \hat{s} >0$, we can consider a nonnegative function $\ell(s)$ such that \mbox{$\ell^2(s)= s^TLs$}.  Taking the gradient of both sides of the latter equation yields $\ell(s) \nabla\ell(s)=Ls$, i.e., (\ref{11/29/23.1}), assuming $L$ is nonsingular.  But according to the equation in the penultimate sentence, $\ell(\alpha s) = \alpha\ell(s)$ for $\alpha$ close to 1.  Thus, applying $\ell(\cdot)$ to both sides of (\ref{11/29/23.1}) we obtain (\ref{11/29/23.2}).  
For Euclidean space, $L$ is the identity and we have $s=\ell(s)\nabla \ell(s)$, $\ell\left(\nabla \ell(s)\right) = 1$, $s\in \mathbb{R}^n$.  As an aside, the latter two equations can actually be derived from the four postulates of George Birkhoff's formulation of Euclidean geometry \cite{birkhoff:1932}, providing a new proof of the Pythagorean Theorem \cite{fgreensite:2022b}.  For Minkowski space and its Minkowski norm (or if one prefers, ``norm"), one of the eigenvalues of $L$ has sign opposite to the sign of all other eigenvalues.
So, distinct geometries result depending on the signature of the scalar product $L$ in the quadratic space setting.  

With the above in mind, we may observe that (\ref{5/22/24.5}), (\ref{5/22/24.6}) define a \mbox{``logarithmic} space" counterpart of a quadratic space as defined by (\ref{11/29/23.1}), (\ref{11/29/23.2}).  That is, for a quadratic space resulting from a scalar product $L$, we have a degree-1 positive homogeneous ``norm" $\ell(s)$ with, \[ Ls = \frac{1}{2}\left[\nabla\ell^2(t)\right]\Big|^{t=s} = \ell(s)\left(\nabla\ell(t)\big|^{t=s}\right).\]  This can be compared to our logarithmic space (which requires an algebra), where we have a degree-1 positive homogeneous ``norm" $\ell_{\rm sp}(s)$ with, \[\frac{1}{\|\mathbf{1}\|^2}\,L_{\shortparallel} s = \nabla\left[\log\ell_{\rm sp}(t)\right]\Big|^{t=s^{-1}} = \ell_{\rm sp}(s)\left(\nabla\ell_{\rm sp}(t)\big|^{t=s^{-1}}\right).\]  

So, for a logarithmic space derived from an algebra, there is a point set given by an appropriate unit ball around $\mathbf{1}$  and there is the ``norm" $\ell_{\rm sp}(s)$.  However, there is also a signed ``angle" that can be assigned to a point $s$ analogous to the way a hyperbolic angle is assigned to a point with respect to the unit hyperbola in $\mathbb{R}^2$ (the signed area of a hyperbolic sector) - e.g., as exploited in Lorentzian 4-space geometry.  Thus, analogous to the case of different signatures of scalar products for a quadratic space, distinct geometries result from $L_{\shortparallel}$ taken from different subspaces of the special unital norm family.



\section{Combinatorial issues}
Regarding computation of the invariants identified in Corollary \ref{8/18/23.2}, we can make the following observations.  As seen in the proof of Proposition \ref{9/2/24.1}, computation of $\mathfrak{u}$ in the form of $M_\mathfrak{u}$ is straightforward, as is the accompanying computation of the dimension of $\mathfrak{u}$.   Computation of the largest rank of a member of $\mathfrak{u}$ is also easy, e.g., using the method of minors.  Computation of the smallest rank of a nonzero member of $\mathfrak{u}$ could be problematic, since it might require evaluation of the dimension of a complicated algebraic set (e.g., potentially involving evaluation of the simultaneous vanishing of many minors).  Evaluation of the number of parameters to which ${\rm det}(M_{\mathfrak{u}})$ is sensitive is simply a matter of inspection.  Examples of invariants derived from $V_\mathbb{R}({\rm det}(M_{\mathfrak{u}}))$ include the real dimension of that real hypersurface (for which an efficient algorithm is \mbox{available \cite{bannwarth:2015})}, and the number of irreducible components of $V_\mathbb{R}({\rm det}(M_{\mathfrak{u}}))$.  

 
Based on the use of one or more simultaneously selected computable invariants identified from Corollary \ref{8/18/23.2}, various more-or-less combinatorial questions arise when the invariants (singly or jointly) are considered in combination with the invariant given by algebra dimension.  For example, with respect to a chosen algebra dimension,  

\begin{enumerate} \item How many values of the anti-rotor dimension occur?  \item How many irreducible components of the various $V_\mathbb{R}({\rm det}(M_{\mathfrak{u}_A}))$ can occur?  \item For which algebras is ${\rm det}(M_{\mathfrak{u}_A})$ identically $0$?  \item For each $k\le \text{dim}(\mathfrak{u}_A)$, can $k$ ever be the maximal rank of a member of $\mathfrak{u}_A$?   \end{enumerate}  
Similar questions arise from the additional invariants identified in Corollary \ref{7/25/24.1}.  Quest for the above characterizations of the set of $n$-dimensional algebras can be considered a response to the intractability of the algebra classification problem.  On the other hand (despite some concrete results listed below), it is likely that many of these combinatorial problems will be very difficult.

Concerning the general question of the distribution of $n$-dimensional algebras into distinct ``bins" (i.e., classes) defined by the anti-rotor dimension, according to the proof of Proposition \ref{9/2/24.1}, $\mathfrak{u}$ is derived from an initial set of $\frac{n(n+1)}{2}$ parameters, which are then subjected to  $\frac{n(n-1)}{2}$ constraint equations.  One might then expect to be finally left with $\frac{n(n+1)}{2} - \frac{n(n-1)}{2} = n$ parameters - i.e., $\text{dim}(\mathfrak{u}) = n$ (so there would only be a single nonempty bin).  Of course, this reasoning is faulty.  Consider the analogous problem of finding the symmetric matrices $L$ that satisfy (\ref{4/7/23.1}) except that $s^{-1}$ is replaced with $s$ in that equation.  One quickly discovers that every constraint equation reduces to $0=0$, so that {\it every} symmetric $(n\times n)$-matrix is a solution (the dimension of the solution space is $\frac{n(n+1)}{2}$).  So, if some of the constraint equations for uncurling metrics were subject to  that kind of thing, or if the $\frac{n(n-1)}{2}$ constraints did not constitute a set of $n$ independent constraints, perhaps $\text{dim}(\mathfrak{u})$ could be some number greater than $n$ (e.g., see the end of Section 12).  Conversely, it might occur that a constraint equation eliminates more than one parameter, so that $\text{dim}(\mathfrak{u})$ turns out to be less than $n$ (along these lines, see the proof of Theorem \ref{4/8/22.1}).  On the other hand, we still might guess that $\text{dim}(\mathfrak{u}) = n$ would define the most popular bin.  In fact, for all three isomorphism classes of two-dimensional algebras, Table I of Section 9 indicates that \mbox{$\text{dim}(\mathfrak{u}) = 2$}.  For the three-dimensional algebras, Table II indicates that five of the six algebra isomorphism classes satisfy $\text{dim}(\mathfrak{u}) = 3$, although the sixth satisfies $\text{dim}(\mathfrak{u}) = 2$.  Moreover, as regards the four-dimensional algebras, it turns out that $\text{dim}(\mathfrak{u}_{M_2(\mathbb{R})}) = 1$ (as follows from Theorem \ref{4/8/22.1}). So a precise understanding of the distribution of $n$-dimensional algebras into the bins/classes defined by $\text{dim}(\mathfrak{u})$ remains a goal.  

However, Corollary \ref{8/17/22.1} is a first general result relevant to a different aspect of the distribution of $n$-dimensional algebras into bins/classes defined by the anti-rotor dimension, as it indicates that if the anti-rotor is one-dimensional and has a nonsingular member then the algebra is simple.  

Another way to look at this is as follows.
For a given integer $n\ge 4$, there are uncountably many isomorphism classes of $n$-dimensional algebras.  To deal with the accompanying unwieldy situation as regards classification, it becomes of interest to instead pursue combinatorial questions regarding the above discrete invariants as a means of characterizing the set of algebras as an entity in itself.   For example, consider the characterization defined by the lattice subset, \[ \mathcal{X}_{\text{dim}(\mathfrak{u})} \equiv \big\{\big(\text{dim}(A),\text{dim}(\mathfrak{u}_A)\big): A \text{ is an algebra}\big\} \subset \mathbb{Z}\times\mathbb{Z},\] where $\text{dim}(\cdot)$ indicates dimension.  If computable, $\mathcal{X}_{\text{dim}(\mathfrak{u})}$ would be a depiction of whether the set of algebras with dimension $n$ and anti-rotor dimension $m$ is empty or not, for any $(n,m)\in \mathbb{Z}\times\mathbb{Z}$.  That is, algebra dimension and anti-rotor dimension are treated in combination, and counting is with respect to 0 (empty) and 1 (not empty). Although the possibility of a full determination of $\mathcal{X}_{\text{dim}(\mathfrak{u})}$ is in doubt, one can hope to reveal many of its distinguishing features (some perhaps in an  asymptotic or order of magnitude sense, as $n$ becomes large). The following are some basic facts concerning $\mathcal{X}_{\text{dim}(\mathfrak{u})}$, beginning with the obvious and becoming progressively less so.

  \begin{enumerate} \item If $n<1$ or $m<0$, then $(n,m)\notin \mathcal{X}_{\text{dim}(\mathfrak{u})}$, \item  For $f(n)$ equal to the cardinality of $\{m\in\mathbb{Z}:(n,m)\in \mathcal{X}_{\text{dim}(\mathfrak{u})}\}$, we have $f(n)\le  1+\frac{n(n+1)}{2}$. \item For algebras of dimension less than or equal to three, the only corresponding members of $\mathcal{X}_{\text{dim}(\mathfrak{u})}$ are (1,1), (2,2), (3,2), and (3,3), so $f(1)=1$, $f(2)=1$, $f(3)=2$.  \item  $(n,n)\in \mathcal{X}_{\text{dim}(\mathfrak{u})}$ for all $n\ge 1$, and for any $n>1$ there are at least two algebra isomorphism classes of $n$-dimensional algebras such that their members all have anti-rotor dimension $n$.  \item $(n,0)\notin \mathcal{X}_{\text{dim}(\mathfrak{u})}$, i.e., nontrivial members of $\mathfrak{u}_A$ always exist.  \item  $(n^2,1)\in \mathcal{X}_{\text{dim}(\mathfrak{u})}$, for $n\ge 1$ (and as indicated in (3) above, $(n,1)$ is not necessarily a member of $\mathcal{X}_{\text{dim}(\mathfrak{u})}$). 
\item If $\text{dim}(\mathfrak{u}_A) = 1$, and $\mathfrak{u}_A$ has a nonsingular member, then $A$ is simple.
\end{enumerate}
  
  \medskip
  
As we have indicated, a combinatorial task of some interest in regard to how the set of algebras is structured, would be to find an expression for $f(n)$ (the number of distinct values of anti-rotor dimension that members of the set of $n$-dimensional algebras actually have) - or at least greatly improve on (2) above.   
There are evidently many similar combinatorial questions that can be proposed to additionally characterize the set of algebras, stemming from others of the anti-rotor based invariants (treated singly or jointly).

\section{Elementary functorial aspects and a couple of useful results}
 Theorem \ref{6/6/23.1} indicates that the association of algebras with anti-rotors is functorial.  As a basic example, let {\bf A} be the category each of whose objects is an algebra, where the morphisms are algebra isomorphisms, and let {\bf S} be the category each of whose objects is a finite-dimensional vector space of real symmetric $(n\times n)$-matrices, where the value of $n\ge 1$ depends on the object (i.e., the members of the object vector space consist of different linear combinations of the members of a finite set of linearly independent real symmetric $(n\times n)$-matrices for some fixed value of $n$).  For objects $S_1,S_2$ of {\bf S}, there is a morphism $k$ if there exists a nonsingular real matrix $k$ such that \mbox{$S_1  = k^T  S_2 k$} (so we denote morphisms and isomorphisms by the same symbol).  Obviously, the identity morphisms exist and the morphisms are associative.    Definition \ref{7/21/23.1} specifies a mapping from algebras (the objects in {\bf A}) to anti-rotors (which are objects in {\bf S}), and (keeping in mind Theorem \ref{6/6/23.1}) the equation in the penultimate sentence indicates how a morphism in {\bf A} selects a morphism in {\bf S}.  It is thus easily verified that (\ref{5/23/23.1}) leads to a functor from {\bf A} to {\bf S}.  This immediately results in a class of algebra isomorphism invariants, some of which are listed in  \mbox{Corollary \ref{8/18/23.2}}.
Continuing in that vein, this program can similarly address the algebra epimorphism problem (determining whether two algebras are related by an epimorphism) using a simple modification of the above functor.  

Thus, let ${\bf E}$ be the category each of whose objects is an algebra (i.e., a real unital associative algebra of nonzero finite-dimension), and a morphism $k$ referable to a source $A$ and target $A'$ exists if there is an algebra epimorphism  $k:A\rightarrow A'$ (obviously the identity morphisms exist, and the morphisms are associative).    

We define category {\bf P} as follows. An object is a symmetric $(n\times n)$-matrix where $n$ is any chosen positive integer, and where each matrix entry is a linear combination of the components of an $m$-tuple of distinct real parameters, $m\ge 1$.  Furthermore, this parametrized matrix has the feature that if all of the $m$ parameters are replaced by 0 except for one parameter being replaced by 1, and one forms the collection of all such matrices except that a different one of the $m$ parameters is replaced by $1$ in each case, then the resulting $m$ real matrices are linearly independent.  Thus, for any object of {\bf P}, the implied set of real matrices resulting from all realizations of the associated $m$-tuple of parameters is an $m$-dimensional vector space (a realization of the $m$-tuple of parameters is the $m$-tuple resulting from a choice of real values for the $m$ parameters).  An object of {\bf P} is considered to be the same as the latter vector space, each of whose members is a real matrix, and can be represented by a matrix whose entries are linear combinations of a set of real parameters.   There is a morphism pointing from object $P$ to object $P'$ if there exists a real matrix $w$ such that $ww^T$ is nonsingular and $w^TPw$ is a subspace of $P'$.  We could similarly notate the morphism as $w$, but it turns out that would cause confusion once we introduce the functor.  Therefore, we will instead suggestively notate the morphism as $w^T[\cdot]w$.    The composition of morphism $\begin{CD}  P @>w^T[\cdot]w>> P'\end{CD}$ with morphism $\begin{CD}  P' @>w'^T[\cdot]w'>> P''\end{CD}$ is defined to be the morphism with source $P$ and target $P''$ resulting from $w'w$ (i.e., analogous to how a morphism from $P$ to $P'$ results from real matrix $w$).  This morphism exists since it is easy to show that $(w'w)^TP(w'w)$ is a subspace of $P''$ and $(w'w)(w'w)^T$ is nonsingular. In other words, given the way we are notating the morphisms, \begin{equation}\label{2/22/25.1}(w'^T[\cdot]w')\circ (w^T[\cdot]w) \equiv (w'w)^T[\cdot](w'w).\end{equation} It is easy to verify that composition of morphisms is associative. The identity morphisms obviously exist. Thus, {\bf P} is a category.

We define a mapping $F$ from category ${\bf E}$ to category {\bf P} as follows.  First, for any object $A$ of {\bf E}, $F[
A] \equiv \mathfrak{u}_A$.  Second, given a morphism $k_1$ from $A_1$ to $A_2$, if $k_1[\mathfrak{u}_{A_2}]k_1^T \subset \mathfrak{u}_{A_1}$ and $k_1k_1^T$ is nonsingular, then the morphism $F(k_1)$ pointing from $F[A_2]$ to $F[A_1]$ is taken to be $\begin{CD}  \mathfrak{u}_{A_2} @>k_1^T[\cdot]k_1>> \mathfrak{u}_{A_1}\end{CD}$ .  Otherwise, $F(k_1)$ is undefined.

\begin{theorem}\label{6/13/22.1}  $F: \mathbf{E}\rightarrow \mathbf{P}$ is a contravariant functor. 
 \end{theorem}
 \begin{proof}  The identity morphisms are obviously preserved.
We will next show show that the category $\mathbf{E}$ morphism $ \begin{CD}  A_1 @>k_1>> A_2\end{CD}$ implies the existence of category {\bf P} morphism $ \begin{CD}  F[A_2] @>F(k_1)>> F[A_1]\end{CD}$.  That is, we will show that $k_1^T[\mathfrak{u}_{A_2}]k_1 \subset \mathfrak{u}_{A_1}$ with $k_1k_1^T$ nonsingular.

 Taking $s_1$ to be notation for a variable in $A_1$, it follows from the epimorphism assumption that $(k_1s_1)^{-1} = k_1(s_1^{-1})$.  By definition, given the anti-rotor $F[A_2]$ associated with $A_2$, and an arbitrary real matrix  $L_2 \in F[A_2]$, we have $d\left((s_2^{-1})^TL_2 (\mathfrak{d}s_2)\right) = 0$ in some neighborhood of ${\bf 1}_{A_2}$, where $s_2$ is notation for a variable in $A_2$.  Since $k_1$ is an epimorphism, we are entitled to specify $s_2=k_1 s_1$ with respect to appropriate neighborhoods of $\mathbf{1}_{A_1}$ and $\mathbf{1}_{A_2}$.   It follows that \begin{eqnarray}\nonumber 0 = d\left((s_2^{-1})^TL_2 \mathfrak{d}s_2\right) &=&  d\left(((k_1s_1)^{-1})^TL_2 \mathfrak{d}(k_1s_1)\right) \\  &=& d\left((s_1^{-1})^T \big[k_1^T L_2k_1\big]\mathfrak{d}s_1\right), \label{7/4/22.5}  \end{eqnarray}   using the equation in the first sentence of this paragraph (the epimorphism assumption), and the fact that a differential commutes with a linear transformation.  Equation (\ref{7/4/22.5}) is simply an instance of (\ref{4/7/23.1}), so it follows that symmetric $k_1^T L_2k_1$ is an uncurling metric associated with $A_1$.  
 Consequently, $k_1^TF[A_2]k_1 = k_1^T\mathfrak{u}_{A_2}k_1$ is a subspace of $F[A_1] = \mathfrak{u}_{A_1}$.  Furthermore, $k_1k_1^T$ is nonsingular since $k_1$ is an epimorphism.  Hence, $k_1^T[\cdot]k_1$ is a morphism pointing from $F[A_2]= \mathfrak{u}_{A_2}$ to $F[A_1]=\mathfrak{u}_{A_1}$.  That is, $F(k_1)=k_1^T[\cdot]k_1$ is well defined.

Finally, consider
      \[  \begin{CD}  
A_1 @>k_1>> A_2 @>k_2>> A_3\\
 @VVV   @VVV  @VVV\\
\mathfrak{u}_{A_1}   @<k_1^T[\cdot]k_1<<  \mathfrak{u}_{A_2}  @<k_2^T[\cdot]k_1<<  \mathfrak{u}_{A_3} \\ 
 \end{CD} \]
Based on the prior paragraph, morphism $F(k_1) = k_1^T[\cdot]k_1$ exists and morphism $F(k_2) = k_2^T[\cdot]k_2$ exists.  But by the definition of how morphisms in $\mathbf{P}$ are composed, i.e., (\ref{2/22/25.1}), we have $ F(k_2\circ k_1) = F(k_1)\circ F(k_2)$.   The theorem then follows.
\end{proof}

This leads to some useful results.  \enlargethispage{1cm}

\begin{corollary} \label{8/17/22.5} If $k_1:A_1\rightarrow A_2$ is an epimorphism, then ${\rm dim}(\mathfrak{u}_{A_2}) \le {\rm dim}(\mathfrak{u}_{A_1})$. \end{corollary}  

\begin{proof} By Theorem \ref{6/13/22.1} and the definition of a morphism in $\mathbf{P}$, we have \mbox{$k_1^T \mathfrak{u}_{A_2} k_1 \subset \mathfrak{u}_{A_1}$} with $k_1k_1^T$ nonsingular.  It follows that  $k_1k_1^T \mathfrak{u}_{A_2} k_1k_1^T \subset k_1\,\mathfrak{u}_{A_1}k_1^T$.  Since $k_1k_1^T$ is nonsingular, we then have 
\begin{equation}\label{2/23/25.1} \mathfrak{u}_{A_2} \subset (k_1k_1^T)^{-1}k_1\, \mathfrak{u}_{A_1} k_1^T(k_1k_1^T)^{-1}.\end{equation}  However, \begin{equation}\label{2/23/25.2}  \text{dim}\left((k_1k_1^T)^{-1}k_1 \,\mathfrak{u}_{A_1} k_1^T(k_1k_1^T)^{-1}\right) \le \text{dim}(\mathfrak{u}_{A_1}).\end{equation} Combining (\ref{2/23/25.1}) and (\ref{2/23/25.2}), the corollary follows.
   \end{proof}
   
  Corollary \ref{8/17/22.5} is useful in excluding the possibility that two algebras are related by an epimorphism (e.g., in excluding an epimorphism from an algebra to one of its subalgebras without having to identify ideals).

The following uses Corollary \ref{1/1/23.1}, which states that the dimension of an algebra's anti-rotor is at least 1. 

\begin{corollary}\label{8/17/22.1} If the anti-rotor of an algebra is one-dimensional and contains a nonsingular member, then the algebra is simple. \end{corollary}

\begin{proof} Suppose algebra $A_1$ is $j$-dimensional and has a non-trivial ideal, and its anti-rotor is one-dimensional and has a nonsingular member.  It follows that we can write $\mathfrak{u}_{A_1} = F[A_1] = \{\alpha L_1:\alpha\in \mathbb{R}\}$, where $L_1$ is a some real nonsingular $(j\times j)$-matrix.  Since $A_1$ has a non-trivial ideal, there exists an algebra epimorphism $k_1:A_1\rightarrow A_2$ where $A_2$ has dimension $q>0$, and $k_1$ has a nontrivial kernel.  Together, \mbox{Corollary \ref{8/17/22.5}} and \mbox{Corollary \ref{1/1/23.1}} indicate that the dimension of the anti-rotor of $A_2$ is 1.   Thus we can write $\mathfrak{u}_{A_2} = F[A_2] = \{\alpha L_2:\alpha\in \mathbb{R}\}$, where $L_2$ is a nonzero real $(q\times q)$-matrix.  It follows that $k_1^T L_2 k_1$ also cannot the zero matrix - since if it were, $k_1k_1^T L_2 k_1k_1^T$ would be the zero matrix, and applying $(k_1k_1^T)^{-1}$ to both sides of the prior expression would indicate that $L_2$ {\it is} the zero matrix.  

From (\ref{7/4/22.5}), we know that $ k_1^T L_2 k_1$ must be a member of $F[A_1]$, and from the prior paragraph we know that $ k_1^T L_2 k_1$ is not the zero matrix.  But the determinant of $k_1^T L_2 k_1$ is zero, since $k_1$ has a  nontrivial kernel.  And yet, the only member of the anti-rotor $F[A_1]$ having zero determinant is the zero matrix (corresponding to $\alpha=0$), since $L_1$ is nonsingular by assumption.   Thus, $k_1$ cannot exist, and so $A_1$ is simple. \end{proof}

In view of Theorem \ref{4/8/22.1}, a consequence of Corollary \ref{8/17/22.1} is an alternative demonstration that the algebra of real $(n\times n)$-matrices $ M_n(\mathbb{R})$ is simple.  

Note that the converse of Corollary \ref{8/17/22.1} is not true.  For example, $\mathbb{C}$ is simple but has a two-dimensional anti-rotor (see Section 9). 

\section{Cohomology and the third family of isomorphism invariants derived from the anti-rotor} 
If the Hochschild cohomology dimension is finite, then a third family of invariants for algebra $A$ results from the anti-rotor of the (real unital associative graded commutative finite-dimensional) algebra derived from the Hochchild \mbox{cohomology} groups via the cup product.  These invariants are descriptive of not only $A$ but also the cohomology ring.  The alternative case of infinite Hochschild cohomology dimension is one motivation for attempting to extend the anti-rotor concept to infinite-dimensional algebras. 


Although in Section 6 there was an identified functor relevant to limited purpose categories relating to some aspects of algebras and anti-rotors, one can quickly dismiss the possibility of a more capable functor between respective abelian categories.  Specifically, the category of $A$-bimodules (supporting the Hochschild cohomology), and a prospective abelian ``anti-rotor category" (i.e., at least encompassing anti-rotors), cannot be associated with a functor from one to the other.  That is, while the objects of the category of $A$-bimodules could each be fitted with the standard product of elements employed in a tensor product of algebras, and thereby be associated with anti-rotors, the bimodule homomorphisms would have nothing at all to do with such entities (i.e, would not respect the imposed product operation), and thus could not be assigned counterparts in the anti-rotor category.  In other words, a chain complex deriving from a chosen algebra in a supposed abelian anti-rotor category would have to stand on its own - though such a prospect is not necessarily excluded.

\section{The dual of the anti-rotor generalizes the usual algebra norm} 

The ``usual" algebra norm is given by the determinant of an element's image under the left regular representation (e.g., \cite{bourbaki:1989}).  The dual of the anti-rotor (as the group of unital norms) subsumes this definition in the sense of the theorems and corollaries of this section (particularly when combined with the observations made toward the end of Section 4).  This also means that the anti-rotor is at least one-dimensional.

 We begin with the matrix algebra $ M_n(\mathbb{R})$. For an invertible element $s$, with det$(s)$ being its determinant, we have the well-known (though in non-standard format),
  \begin{equation}\label{11/5/21.1} \nabla \mbox{det}(s) = \big[{\rm Adj}(s)\big]^T = \mbox{det}(s)[s^{-1}]^T,\end{equation} where Adj$(s)$ is the adjugate matrix.  The first of these equalities is special case of Jacobi's formula, and the second follows from the classical matrix inverse expression resulting from Laplace expansion.  The gradient notation on the left-hand-side is employed with the understanding that an element of the algebra (including the full expression on the right-hand-side), as a member of the vector space of elements, is identified as the $n^2$-tuple obtained from sequentially appending the columns of its matrix format (and vice-versa), with a compatible understanding  regarding the transpose operation as a particular involutive $(n^2\times n^2)$-matrix, which we will denote as $\mathcal{T}$. 
  
From (\ref{11/5/21.1}), $\mathcal{T}s^{-1} = \frac{\nabla{\rm det}(s)}{{\rm det}(s)} = \nabla[\log{\rm det}(s)]$, so $\mathcal{T}$ is an uncurling metric.  Since ${\rm det}(s)$ is a degree-$n$ positive homogeneous function, the Euler Homogeneous Function Theorem implies $[\nabla{\rm det}(s)]\cdot s = n \,{\rm det}(s)$.  We also have $\|\mathbf{1}\|^2 = n$.  Consequently, (\ref{7/7/21.1}) is satisfied for $L=\mathcal{T}$, so $\mathcal{T}$ is a normalized uncurling metric.   Applying the first equation of this paragraph to (\ref{11/7/20.2}), it is then seen that $(\text{det}(s))^{\frac{1}{n}}$ is a special unital norm.
  Moreover,
 
 \begin{theorem} \label{4/8/22.1} $\mathfrak{u}_{M_n(\mathbb{R})}$ is one-dimensional and there is only a single normalized uncurling metric, this being the  transpose operation on $(n\times n)$-matrices expressed as an $(n^2\times n^2)$-matrix. The associated special unital norm is $({\rm det}(s))^{\frac{1}{n}}$, i.e., the usual norm raised to the power given by the reciprocal of the dimension of $M_n(\mathbb{R})$. \end{theorem}
 \begin{proof} The vector space of members of $ M_n(\mathbb{R})$ consists of  vectors with $n^2$ components.   Using the sequentially appended matrix columns convention as introduced in the second paragraph of this section, we can equivalently represent a member $s\in M_n(\mathbb{R})$ in either matrix or column vector format.   An uncurling metric for this algebra will be a matrix $L\in M_{n^2}(\mathbb{R})$ that must satisfy (\ref{4/7/23.1}), i.e., $d\left([Ls^{-1}]\cdot \mathfrak{d}s\right)=0$. This constitutes a condition on every pair of distinct components of $Ls^{-1}$ (the latter again having $n^2$ components).  Thus,  the partial derivative of the  $(i,j)$-entry of matrix $Ls^{-1}$ (say), taken with respect to the $(k,m)$-entry of matrix $s$ (say) (with $(k,m)\ne (i,j)$),  minus  the partial derivative of the $(k,m)$-entry of $Ls^{-1}$  taken with respect to the $(i,j)$-entry of matrix $s$, must be zero.
 
 To examine this constraint, let $s_{pq}$ be the $(p,q)$-entry of matrix $s$.   Differentiating both sides of the expression
$ ss^{-1} = {\bf 1}$ with respect to $s_{pq}$ (as in Matrix Calculus), we obtain $\frac{\partial (s^{-1}s)}{\partial s_{pq}} = \frac{\partial (s^{-1})}{\partial s_{pq}}s + s^{-1}\frac{\partial s}{\partial s_{pq}}=0$ (where $0$ is the matrix of all zero entries).  Thus, we have the well-known result 
 \[ \frac{\partial (s^{-1})}{\partial s_{pq}} = -s^{-1}\frac{\partial s}{\partial s_{pq}}s^{-1} = -s^{-1}e_{pq}s^{-1},\]   where $e_{pq}$ is the matrix whose $(p,q)$ entry is $1$, but with all other entries zero (and $\{ e_{pq}\}$ is the canonical basis for $ M_n(\mathbb{R})$).  
 Hence, the partial derivative of $s^{-1}$ with respect to the  $(k,m)$-entry of matrix $s$, which is itself a matrix, (``the derivative of $s^{-1}$ in the $e_{km}$-direction") is, \begin{equation}\label{8/18/22.5} -s^{-1} e_{km}s^{-1} = -(s^{-1})_{:k} (s^{-1})_{m:},\end{equation} where $(s^{-1})_{:k}$ indicates the $k$-th column of $s^{-1}$, and $(s^{-1})_{m:}$ indicates the $m$-th row of $s^{-1}$.    Similarly, the derivative of $s^{-1}$ with respect to the $(i,j)$-entry of $s$ (``the derivative of $s^{-1}$ in the $e_{ij}$-direction")  is, \begin{equation}\label{8/18/22.6} -s^{-1} e_{ij}s^{-1} = -(s^{-1})_{:i} (s^{-1})_{j:}.\end{equation}  It is thus clear that the sets of entries of the $e_{km}$-direction derivative matrix and the $e_{ij}$-direction derivative  matrix share the member  $-(s^{-1})_{mi} (s^{-1})_{jk}$, and they can share no other member because $\{(s^{-1})_{pq}\} = \{\frac{C_{qp}}{\det(s)}\}$, where the right-hand-side of this equation is clearly a set of linearly independent functions since $\{C_{qp}\}$ is the matrix of cofactors of $s$.  The latter observation also indicates that the union of the entries of the $e_{km}$-direction derivative matrix with the entries of the $e_{ij}$-direction derivative matrix is a set of linearly independent functions.

Therefore, since a derivative and a linear transformation commute, the partial derivative in the $e_{km}$-direction of the $e_{ij}$-component of $Ls^{-1}$ is a linear combination of the $n^2$ entries of the matrix on the right-hand-side of (\ref{8/18/22.5}), and  the partial derivative in the $e_{ij}$-direction of the $e_{km}$-component of $Ls^{-1}$ is a linear combination of the $n^2$ entries on the right-hand-side of (\ref{8/18/22.6}).
So satisfaction of \mbox{$d\left([Ls^{-1}]\cdot \mathfrak{d}s\right) = 0$,} based on the observation made at the end of the prior paragraph, requires that all entries in any two rows of $L$ must be zero, except for the one entry in each row corresponding to the shared member of the two relevant partial derivative matrices (the right-hand-sides of (\ref{8/18/22.5}) and (\ref{8/18/22.6})) noted in the prior paragraph - and these two corresponding nonzero entries of the two cited rows of $L$ must have the same value.  This prescription must be true for any two rows of $L$, so all nonzero entries of $L$ must have the same value, and there can be at most one nonzero value in each row of $L$.  Thus, the anti-rotor is at most one-dimensional. Assuming the group is one-dimensional, the anti-rotor can be represented as a parameter $\alpha$ ranging over $\mathbb{R}$ multiplying a nonzero constant matrix whose nonzero entries are all 1.  
For the special unital norm, the value of $\alpha$ is determined by the requirement that (\ref{7/7/21.1}) be satisfied. 


Hence, there is at most one possible normalized uncurling metric.  But according to observations made in the third paragraph of this section, $ M_n(\mathbb{R})$ does indeed have a normalized uncurling metric - this being the $(n^2\times n^2)$-matrix corresponding to the matrix transpose operation on $(n\times n)$-matrices.  As was shown immediately prior to the theorem statement, $(\mbox{det}(s))^{\frac{1}{n}}$ is the associated special unital norm.  Since the usual norm of an element $s\in M_n(\mathbb{R})$ is $(\mbox{det}(s))^n$, the theorem follows.
 \end{proof}
 
 Expanding on Theorem \ref{4/8/22.1}, the following lemma and theorem indicate that the unital norms are an elaboration of the usual algebra norm.  Recall that a subalgebra of $M_n(\mathbb{R})$ is defined to be {\it self-adjoint} if the transpose operation maps the subalgebra to itself.
 
  \begin{lemma}\label{5/13/22.1} If the left regular representation of an $n$-dimensional algebra $\bar{A}$ with elements $\bar{s}$ is self-adjoint as a subalgebra of $M_n(\mathbb{R})$, then there is a special unital norm given by the usual algebra norm raised to the power of $\frac{1}{n}$, i.e., $\ell_{\rm sp}(\bar{s}) = [{\rm det}(\mathcal{I}\, \bar{s})]^\frac{1}{n}$, where $\mathcal{I}$ is the isomorphism from $\bar{A}$ to its left regular representation.
\end{lemma}

\begin{proof} 
Let $\bar{A}$ be an $n$-dimensional algebra and let $A$ be its left regular representation.  Given $\bar{s}\in\bar{A}$ we denote its image under the left regular representation as $s = \mathcal{I}\bar{s}$.
 So, $\bar{A}$ is isomorphic to $A$, and $A$ is a subalgebra of $ M_n(\mathbb{R})$.   $\mathbf{1}_{M_n(\mathbb{R})}$ is the $(n\times n)$-identity matrix, and observe that as an element of $ M_n(\mathbb{R})$, $\mathbf{1}_A$ is likewise the $(n\times n)$-identity matrix.  However, $ M_n(\mathbb{R})$ is $n^2$-dimensional with $\|\mathbf{1}_{M_n(\mathbb{R})}\|^2 = n$, while $A$ is $n$-dimensional with $\|\mathbf{1}_A\|^2 = \|\mathbf{1}_{\bar{A}}\|^2$ (because $\mathbf{1}_{\bar{A}}$ and  $\mathbf{1}_A$ can be represented by the same $n$-tuple in the standard basis).   

As noted shortly prior to the statement of Theorem \ref{4/8/22.1}, the transpose operation on $ M_n(\mathbb{R})$ is a linear transformation $\mathcal{\mathcal{T}}:\mathbb{R}^{n^2}\rightarrow \mathbb{R}^{n^2}$ and, also notating it as $\mathcal{\mathcal{T}}$ with respect to the standard basis, matrix $\mathcal{\mathcal{T}}$ represents a normalized uncurling metric associated with algebra  $ M_n(\mathbb{R})$.  Being an uncurling metric means that integration of $\mathcal{\mathcal{T}}s^{-1}$ is path-independent in some open ball centered at the multiplicative identity of $ M_n(\mathbb{R})$.  
 Hence, $\int_{\mathbf{1}_{M_n({\mathbb{R})}}}^s (\mathcal{\mathcal{T}}t^{-1})\cdot \mathfrak{d}t$ is path-independent in the latter neighborhood and, in particular, is path-independent for paths in that neighborhood restricted to be in $A$.  By assumption, $\mathcal{\mathcal{T}}$ maps $A$ to $A$. 
 Thus, $\int_\mathbf{1}^s \mathcal{\mathcal{T}}t^{-1}\cdot \mathfrak{d}t$ is path-independent when the domain of integration is restricted to $A$, and on such a path the integrand also remains in $A$.  This means that 
  $d((\mathcal{\mathcal{T}}s^{-1})\cdot \mathfrak{d}s) = 0$ in some neighborhood of $\mathbf{1}_A$ in $A$.
  
   Now, we can use $\mathcal{\mathcal{T}}$ to define a scalar product on $A$ (not necessarily positive-definite or nondegenerate) as $\langle s_1,s_2\rangle_A \equiv  \frac{\|\mathbf{1}_{A}\|^2}{n}s_1^T \mathcal{\mathcal{T}} s_2$. 
    Since $A$ and $\bar{A}$ are isomorphic, the above scalar product on $A$ can be employed to produce a scalar product on $\bar{A}$ so that the two resulting quadratic spaces are isomorphic (i.e., with scalar products that are invariant with respect to the application of $\mathcal{I}$).  Specifically, for $\bar{s_1},\bar{s_2}\in\bar{A}$, define $s_1 \equiv \mathcal{I}\bar{s}_1$ and $s_2 \equiv \mathcal{I}\bar{s}_2$.  Then define the induced scalar product on $\bar{A}$ via, \begin{equation}\label{7/15/22.1} \bar{\mathcal{\mathcal{T}}} \equiv  \frac{\|\mathbf{1}_{\bar{A}}\|^2}{n}\mathcal{I}^T\mathcal{\mathcal{T}}\mathcal{I},\end{equation} so that $\langle \bar{s}_1,\bar{s}_2\rangle_{\bar{A}} \equiv \bar{s}_1^T \bar{\mathcal{\mathcal{T}}} \bar{s}_2$.  Applying (\ref{7/15/22.1}) and the previously noted $\|\mathbf{1}_A\|^2 = \|\mathbf{1}_{\bar{A}}\|^2$, we have \begin{equation}\label{8/25/22.1} \langle \bar{s}_1,\bar{s}_2\rangle_{\bar{A}} = \bar{s}_1^T \bar{\mathcal{T}} \bar{s}_2 = \frac{\|\mathbf{1}_{\bar{A}}\|^2}{n}( s_1^T \mathcal{T} s_2) = \langle s_1,s_2\rangle_{A}.\end{equation}  The fraction $\frac{\|\mathbf{1}_{\bar{A}}\|^2}{n}$ is employed in (\ref{7/15/22.1}) in order that $\langle \mathbf{1}_{\bar{A}},\mathbf{1}_{\bar{A}}\rangle_{\bar{A}} = \|\mathbf{1}_{\bar{A}}\|^2$, given that $\mathbf{1}_{M_n(\mathbb{R})}^T \mathcal{T}\,\mathbf{1}_{M_n(\mathbb{R})} = n = \|\mathbf{1}_{M_n(\mathbb{R})}\|^2$. Due to this isomorphism of $A$ and $\bar{A}$ as quadratic spaces, the previously demonstrated $d((\mathcal{T}s^{-1})\cdot \mathfrak{d}s) = 0$ in some neighborhood of $\mathbf{1}_A$ in $A$ implies 
\mbox{$d((\bar{\mathcal{T}}(\bar{s})^{-1})\cdot d\bar{s}) = 0$.} Thus, $\bar{\mathcal{T}}$ is an uncurling metric for $\bar{A}$.


In the context of $M_n(\mathbb{R})$, we know from Theorem \ref{4/8/22.1} that $L=\mathcal{T}$ satisfies (\ref{7/7/21.1}), where $\|\mathbf{1}\|^2$ in  (\ref{7/7/21.1}) is $\|\mathbf{1}_{M_n({\mathbb{R})}}\|^2 = n$.  Combining these facts with (\ref{8/25/22.1}), we have \begin{equation}\label{9/27/23.1} (\bar{s})^T \bar{\mathcal{T}} (\bar{s})^{-1} = \frac{\|\mathbf{1}_{\bar{A}}\|^2}{n}( s^T \mathcal{T} s^{-1}) = \|\mathbf{1}_{\bar{A}}\|^2 .\end{equation} 
That is, $\mathcal{\bar{T}}$ is a normalized uncurling metric for $\bar{A}$.  Hence, a special unital norm on $\bar{A}$ results from (\ref{11/7/20.2}) (with $\mathbf{1}=\mathbf{1}_{\bar{A}}$)  via the normalized uncurling metric $L_{\shortparallel} = \bar{\mathcal{T}}$.  
  We then obtain, 
\begin{eqnarray}\nonumber \log\ell_{\rm sp}(\bar{s}) = \frac{1}{\|\mathbf{1}_{\bar{A}}\|^2} \int_{\mathbf{1}_{\bar{A}}}^{\bar{s}} (\bar{\mathcal{T}}(\bar{s})^{-1})\cdot \mathfrak{d}\bar{s}
 &=& \frac{1}{\|\mathbf{1}_{\bar{A}}\|^2} \int_{\mathbf{1}_{A}}^{s} \frac{\|\mathbf{1}_{\bar{A}}\|^2}{n}(\mathcal{T}s^{-1})\cdot \mathfrak{d}s \\ &=& \frac{1}{n} \log\text{det}(s),\label{9/27/23.2}\end{eqnarray} where the first equality is from (\ref{11/7/20.2}), the second equality follows from application of (\ref{7/15/22.1}), and the final equation results from Theorem \ref{4/8/22.1}.
The lemma then follows. 
 \end{proof} 
 

So now we consider the general case.
 
  \begin{theorem}\label{11/17/22.1}
  Given any algebra $\bar{A}$, suppose there is an epimorphism $\Omega:\bar{A}\rightarrow \underline{A}$ where $ \underline{A}$ has dimension $m\ge 1$ and such that the left regular representation of $\underline{A}$ is self-adjoint as a subalgebra of $M_m(\mathbb{R})$.  Then there is a member of the special unital norm family of  $\bar{A}$ given by $[{\rm det}(\underline{\mathcal{I}}\Omega \bar{s})]^{\frac{1}{m}}$, for $\bar{s}\in\bar{A}$, where $\underline{\mathcal{I}}$ is the isomorphism from $\underline{A}$ to its left regular representation.  Furthermore, at least one such epimorphism $\Omega$ exists. 
   \end{theorem}
 
  \begin{proof}  
  If $\Omega$ exists, then by Lemma \ref{5/13/22.1} there exists a normalized uncurling metric $\underline{L}_{\shortparallel}$ for $\underline{A}$, where the associated special unital norm of applicable elements $\underline{s}\in \underline{A}$ is, \begin{equation}\label{1/9/25.1} \ell_{\rm sp}(\underline{s}) =    \exp\left(\frac{1}{\|\mathbf{1}_{\underline{A}}\|^2}\int_{\mathbf{1}_{\underline{A}}}^{\underline{s}} [\underline{L}_{\shortparallel}\underline{t}^{-1}]\cdot \mathfrak{d}\underline{t}\right) = [\text{det}(\underline{\mathcal{I}}\, \underline{s})]^\frac{1}{m}.\end{equation}  It is easy to see that this induces a normalized uncurling metric for $\bar{A}$ given by, 
  \begin{equation}\label{1/9/25.2} \bar{L}_{\shortparallel} \equiv \frac{\|\mathbf{1}_{\bar{A}}\|^2}{\|\mathbf{1}_{\underline{A}}\|^2} \Omega^T \underline{L}_{\shortparallel} \Omega, \end{equation} since (\ref{7/7/21.1}) is satisfied because $\Omega$ is an epimorphism (i.e., $\Omega\bar{s} =  \underline{s}$ implies \mbox{$\Omega(\bar{s})^{-1} =  \underline{s}^{-1}$)}.
  
  According to (\ref{11/7/20.2}), $\bar{L}_{\shortparallel}$ is then associated with the special unital norm on $\bar{A}$ given by $\ell_{\rm sp}(\bar{s}) = [\text{det}(\underline{\mathcal{I}}\Omega\bar{s})]^{\frac{1}{m}}$.  That is, in evaluation of the integral resulting from  (\ref{11/7/20.2}), a path of integration from $\mathbf{1}_{\bar{A}}$ to $\bar{s}$ can be replaced by a path of integration from $\mathbf{1}_{\underline{A}}$ to $\underline{s}$ because of the projection implied by the integrand resulting from (\ref{1/9/25.2}).   Furthermore, on this path,
  \begin{eqnarray}\nonumber  [\bar{L}_{\shortparallel}(\bar{t})^{-1}]\,\cdot\, \mathfrak{d}\bar{t} &=& [(\bar{t})^{-1}]^T \bar{L}_{\shortparallel}\,\mathfrak{d}\bar{t} =  \frac{\|\mathbf{1}_{\bar{A}}\|^2}{\|\mathbf{1}_{\underline{A}}\|^2}[(\bar{t})^{-1}]^T \Omega^T \underline{L}_{\shortparallel} \, \Omega \, \mathfrak{d}\bar{t}  = \frac{\|\mathbf{1}_{\bar{A}}\|^2}{\|\mathbf{1}_{\underline{A}}\|^2}\,[\underline{t}^{-1}]^T\, \underline{L}_{\shortparallel}\,\mathfrak{d}\underline{t} \\ &=& \frac{\|\mathbf{1}_{\bar{A}}\|^2}{\|\mathbf{1}_{\underline{A}}\|^2}\,[\underline{L}_{\shortparallel}\underline{t}^{-1}]\,\cdot\, \mathfrak{d}\underline{t}.\label{2/5/25.2}\end{eqnarray}   In view of (\ref{2/5/25.2}), evaluation of (\ref{11/7/20.2}) then yields the right-hand-side of (\ref{1/9/25.1}). 
  

  It only remains to show that there always exists at least one such $\Omega$. 
  
  Evidently, either $\bar{A}$ is simple, or there exists an epimorphism $\Omega:\bar{A}\rightarrow \underline{A}$ where $ \underline{A}$ is a simple nonzero algebra.   By a famous theorem of Wedderburn, a simple nonzero algebra is isomorphic to one of the real algebras $M_q(\mathbb{R})$, or $M_q(\mathbb{C})$, or $M_q(\mathbb{H})$, for some positive integer $q$.  The left regular representation of each of these algebras is self-adjoint.  Thus, if $\bar{A}$ is simple, the theorem conclusion immediately follows by Lemma \ref{5/13/22.1}.  However, if $\bar{A}$ is not simple, the theorem conclusion follows from the first sentence of this paragraph and the first two paragraphs of this proof.
  \end{proof}

 The following is an alternative statement of part of the above theorem.   
 
 \begin{corollary}\label{1/1/23.1} Anti-rotors always have dimension greater than or equal to 1. \end{corollary}


Half of the examples of associative algebras $\bar{A}$ collected in Table III of \mbox{Section 10} are such that their left regular representations are not self-adjoint as subalgebras of $M_n(\mathbb{R})$, thereby illustrating Theorem \ref{11/17/22.1}.  Calculations related to such algebras are explicitly highlighted in the examples of Section 10.4 and Section 10.5.3.

 

\section{Illustrations in dimensions two and three}
Up to an isomorphism, there are only three two-dimensional algebras.
For the algebra of complex numbers, $\mathbb{C}$, we can write its elements $s$ as $(x,y)$, with the product given by $(x,y)(z,w) = (xz-yw, yz+xw)$, with $\mathbf{1}=(1,0)$, so that $\|\mathbf{1}\|^2=1$.  Also, $(x,y)^{-1} = (\frac{x}{x^2+y^2},\frac{-y}{x^2+y^2})$.  According to the algorithm indicated in the proof of Proposition \ref{9/2/24.1}, we start out by writing  $\mathfrak{u}_\mathbb{C} = \left[\begin{array}{ccc}
 \alpha & \beta   \\ 
 \beta & \gamma    
\end{array}\right]$.  \mbox{We then} obtain  $\mathfrak{u}_\mathbb{C}s^{-1} = \left[\begin{array}{ccc}
 \alpha & \beta   \\ 
 \beta & \gamma    
\end{array}\right]\left(\begin{array}{ccc}
x   \\ 
y   
\end{array}\right)^{-1} = \left(\begin{array}{ccc}
\frac{\alpha x - \beta y}{x^2+y^2}  \\ 
\frac{\beta x - \gamma y}{x^2+y^2} 
\end{array}\right)$.  As regards the right-hand-side of this equation, satisfaction of (\ref{4/7/23.1}) requires that the partial derivative with respect to $x$ of the second component must equal the partial derivative with respect to $y$ of the first component.  Consequently, $\mathfrak{u}_{\mathbb{C}} = \left[\begin{array}{ccc}
 \alpha & \beta   \\ 
 \beta & -\alpha.    
\end{array}\right]$.  Using Definition \ref{9/21/23.2}, a simple exercise in integration leads to $\log  \mathfrak{u}^d_{\mathbb{C}} = \frac{\alpha}{2}\log(x^2+y^2) + \beta{\rm arctan}(\frac{y}{x})$.

The split-complex numbers $\slashed{\mathbb{C}}$ constitute an algebra for which the elements are  $(x,y)$ with product $(x,y)(z,w) = (xz+yw,xw+yz)$, $ \mathbf{1} = (1,0)$, and $\|\mathbf{1}\|^2 = 1$.  Also, $(x,y)^{-1} = (\frac{x}{x^2-y^2},\frac{-y}{x^2-y^2})$. \hspace{1cm}  We then find $\mathfrak{u}_{\slashed{\mathbb{C}}} = \left[\begin{array}{ccc}
 \alpha & \beta   \\ 
 \beta & \alpha.    
\end{array}\right]$, \hspace{.5cm} and \hfill \\ $\log  \mathfrak{u}^d_{\slashed{\mathbb{C}}} = \frac{\alpha}{2}\log[(x+
y)(x-y)] +  \frac{\beta}{2}\log \frac{x+y}{x-y}$.  In fact, this can be simplified to $\log  \mathfrak{u}^d_{\slashed{\mathbb{C}}} = \frac{\alpha+\beta}{2}\log(x+y) + \frac{\alpha-\beta}{2}\log(x-y)$.  Of course, $\frac{\alpha+\beta}{2}$, $\frac{\alpha-\beta}{2}$ are themselves two independent parameters that we could name $\alpha'$, $\beta'$.

The final real two-dimensional unital algebra (up to an isomorphism) is the algebra of dual numbers $\mathbb{D}$.  The elements are $(x,y)$ with product given by $(x,y)(z,w) = (xz,xw+yz)$, with $\mathbf{1}=(1,0)$ and $\|\mathbf{1}\|^2 = 1$. Also, $(x,y)^{-1} = (\frac{1}{x},-\frac{y}{x^2})$.  We then find $\mathfrak{u}_\mathbb{D} = \left[\begin{array}{ccc}
 \alpha & \beta   \\ 
 \beta & 0    
\end{array}\right]$ and $\log  \mathfrak{u}^d_{\mathbb{D}} = \alpha\log x + \beta\left(\frac{y}{x}\right)$. 

   
   Now we consider $\mathbb{R}\times\mathbb{R}$ (understood to be a direct product of rings), which is isomorphic to $\slashed{\mathbb{C}}$.  We have elements as $(x,y)$, product $(x,y)(z,w) = (xz,yw)$, $\mathbf{1}=(1,1)$, and $\|\mathbf{1}\|^2 = 2$.  Also $(x,y)^{-1} = (\frac{1}{x},\frac{1}{y})$.  We then find $\mathfrak{u}_{\mathbb{R}\times\mathbb{R}} = \left[\begin{array}{ccc}
 \alpha & 0   \\ 
 0 & \gamma    
\end{array}\right]$ and $\log  \mathfrak{u}^d_{\mathbb{R}\times\mathbb{R}} = \frac{\alpha}{2}\log x + \frac{\gamma}{2}\log y$. 



With the observations preceding Corollary \ref{7/25/24.1} in mind, let us look further at the isomorphic algebras $\slashed{\mathbb{C}}$ and $\mathbb{R}\times\mathbb{R}$.  It is clear that $\log\mathfrak{u}^d_{\slashed{\mathbb{C}}}$ and $\log\mathfrak{u}^d_{\mathbb{R}\times\mathbb{R}}$ as calulated above are consistent based on the algebra isomorphism resulting from $\left[\begin{array}{ccc}
 1 & -1   \\ 
 1 & 1    
\end{array}\right]\left(\begin{array}{ccc}
x   \\ 
y   
\end{array}\right) = \left(\begin{array}{ccc}
x-y   \\ 
x+y   
\end{array}\right)$.  However, if we ignore that transformation, we could alternatively calculate (\ref{7/29/24.1}), which requires evaluations only on the respective coordinate axes of $\slashed{\mathbb{C}}$ and $\mathbb{R}\times\mathbb{R}$ as a means of assessing the possibility of an isomorphism between the latter two algebras (as recommended prior to the statement of Corollary \ref{7/25/24.1}).  Based on the above calculations, $\slashed{\mathbb{C}}$ results in (\ref{7/29/24.1}) as $\alpha\log x+ \frac{\alpha}{2}\log[(1+
y)(1-y)] +  \frac{\beta}{2}\log \frac{1+y}{1-y}$ and $\mathbb{R}\times\mathbb{R}$ results in (\ref{7/29/24.1}) as $\frac{\alpha}{2}\log x +\frac{\beta}{2} \log y$. Consequently,  \[(\tau_{\slashed{\mathbb{C}},{\rm rat}},\tau_{\slashed{\mathbb{C}},{\rm log}},\tau_{\slashed{\mathbb{C}},{\rm arc}}) = (\tau_{\mathbb{R}\times\mathbb{R},{\rm rat}},\tau_{\mathbb{R}\times\mathbb{R},{\rm log}},\tau_{\mathbb{R}\times\mathbb{R},{\rm arc}}) = (0,2,0). \] Thus, the invariants in Corollary \ref{7/25/24.1} are at least consistent with $\slashed{\mathbb{C}}$ and $\mathbb{R}\times\mathbb{R}$ being isomorphic.  Similarly applying Corollary \ref{7/25/24.1}, it is easily seen that neither of $\slashed{\mathbb{C}}$ or $\mathbb{R}\times\mathbb{R}$ can be isomorphic to either $\mathbb{C}$ or $\mathbb{D}$ (nor can  $\mathbb{C}$ or $\mathbb{D}$ themselves be isomorphic).   So, even though it is in general impractical to search for an isomorphism mapping for algebras that might be isomorphic, invariants derived from $\log\mathfrak{u}^d$ are available to coarsely sort the algebras in a meaningful way.

Incidentally, if we multiply $\mathfrak{u}_{\mathbb{C}}$ on the left by the exchange matrix $\left[\begin{array}{ccc}
 0 & 1   \\ 
 1 & 0    
\end{array}\right]$, and also multiply $\mathfrak{u}_{\mathbb{D}}$ on the left by $\left[\begin{array}{ccc}
 0 & 1   \\ 
 1 & 0    
\end{array}\right]$, we obtain the left regular representations of $\mathbb{C}$ and $\mathbb{D}$, respectively.   In fact, $\mathfrak{u}_{\slashed{\mathbb{C}}}$ and $\mathfrak{u}_{\mathbb{R}\times\mathbb{R}}$ are already the same as the left regular representations of $\slashed{\mathbb{C}}$ and $\mathbb{R}\times\mathbb{R}$, respectively.  This type of phenomenon is {\it not} typical in higher dimensions, apart from special cases like the algebras given by a direct product $\prod_{i=1}^k \mathbb{R}$, and the upper triangular Toeplitz matrix algebras.

As is evident from Theorem \ref{11/17/22.1}, all unital norm groups have a one dimensional subgroup corresponding to multiples of a function closely allied with the usual algebra norm, and the parameter of $M_{\mathfrak{u}_A}$ corresponding to this subgroup {\it always} makes a contribution to $\tau_{A,\log}$.  Thus, in distinguishing algebras, it makes sense to consider a ``reduced triple" of the invariants in Corollary \ref{7/25/24.1} as $(\tau_{A,{\rm rat}},\tau_{A,{\rm log}}-1,\tau_{A,{\rm arc}})$. 
 The three isomorphism classes of the two-dimensional algebras represented by $\mathbb{C}$, $\slashed{\mathbb{C}}$, and $\mathbb{D}$, respectively correspond to the reduced triples  $(0,0,1)$, $(0,1,0)$, and $(1,0,0)$ (i.e., pure elemental forms of the reduced triples associated with algebras in general).  
 
 
 
 So \dots,
 
 \begin{itemize}
\item There is only a single isomorphism class for the one-dimensional algebras, a trivial case in all respects (its reduced triple is (0,0,0)).
\item
As we have noted, there are three isomorphism classes of two-dimensional algebras (with the reduced triples respectively as (0,0,1), (0,1,0), and (1,0,0)). These are listed in Table I, along with their left regular representations, anti-rotors, and unital norm groups. 
The six isomorphism classes of three-dimensional algebras are also instructive.  The structure constants for a representative of each of these are presented in \cite{kobayashi:2021}.  Table II lists those six algebras (in the order given in \cite{kobayashi:2021}) in matrix representations along with their respective anti-rotors and unital norm groups. 
\item 
In contrast, for a choice of $n>3$, there are uncountably many isomorphism classes for the set of $n$-dimensional algebras.
\end{itemize}

\bigskip

Despite the many obvious means for classification of the $n$-dimensional algebras using anti-rotor based invariants for $n\le 3$, one is also entitled to characterize combinatorially in these low dimensions.  Indeed, there are many possible combinatorial questions related to the identified invariants that can be examined, singly or jointly.  As we have observed earlier, five of the six three-dimensional algebra isomorphism classes are associated with an anti-rotor dimension equal to the algebra dimension, whereas for the two-dimensional algebras all isomorphism classes have an associated anti-rotor dimension equal to the algebra dimension.  This begs the question of the combinatorial relationship of algebra dimension and anti-rotor dimension for algebras of dimension greater than three (e.g., for a choice of $n$, how many different anti-rotor dimensions occur for the set of $n$-dimensional algebras?).  Apart from the potential interest of such individual combinatorial questions, classification in higher dimensions does not appear to be possible, while such combinatorial characterizations may be.
 
\vspace{2cm}

  \centerline{\Large The Three Isomorphism Classes of the Two-dimensional Algebras } 

\bigskip


\hspace{.7cm}$ \begin{array}{cccccccc} 
\mbox{\underline{algebra}} & \mbox{\underline{anti-rotor $(M_{\mathfrak{u}})$}} & \mbox{\underline{unital norm group ($\mathfrak{u}^d$)}}  \\ 
\\
\mathbb{C} \approx \left[\begin{array}{ccc}
x  & -y    \\
 y & x \\   
\end{array}\right]   & \left[\begin{array}{ccc}
\alpha  & \beta    \\
 \beta & -\alpha \\   
\end{array}\right]  & (x^2+y^2)^{\frac{\alpha}{2}} \,e^{\beta\,{\rm arctan}\left(\frac{y}{x}\right)}  \\ \\
\slashed{\mathbb{C}} \approx \left[\begin{array}{ccc}
x  & y    \\
 y & x \\   
\end{array}\right]   & \left[\begin{array}{ccc}
\alpha  & \beta    \\
 \beta & \alpha \\   
\end{array}\right]  & [(x+y)(x-y)]^{\frac{\alpha}{2}} \left[\frac{x+y}{x-y}\right]^{\frac{\beta}{2}}  \\ \\ 
\mathbb{R}\times\mathbb{R} \approx \left[\begin{array}{ccc}
x  & 0    \\
 0 & y \\   
\end{array}\right]   & \left[\begin{array}{ccc}
\alpha  & 0    \\
 0 & \beta \\   
\end{array}\right]  & x^\frac{\alpha}{2} y^\frac{\beta}{2}  \\ \\
\mathbb{D} \approx \left[\begin{array}{ccc}
x  & y    \\
 0 & x \\   
\end{array}\right]   & \left[\begin{array}{ccc}
\alpha  & \beta    \\
 \beta & 0 \\   
\end{array}\right]  & x^\alpha\,e^{\beta \frac{y}{x}}  \\ \\
\end{array} $

\bigskip


\noindent {\bf Table I.} {\small Each row indicates a representative  of one of the three two-dimensional algebra isomorphism classes along with its left regular representation, followed by the associated anti-rotor/uncurling metric group and unital norm group (``$\approx$" means \mbox{``isomorphic to"}). Note that  $\slashed{\mathbb{C}}$ and $\mathbb{R}\times\mathbb{R}$ are isomorphic, their respective uncurling metric groups are compatible (via the algebra isomorphism $\left[\begin{array}{ccc}
 1 & 1   \\ 
 -1 & 1    
\end{array}\right]$), and their unital norm groups are  compatible (i.e., $(x+y)^{\frac{\alpha+\beta}{2}}(x-y)^{\frac{\alpha-\beta}{2}}$ versus $x^\frac{\alpha}{2} y^\frac{\beta}{2}$).  }

\pagebreak

\hspace{0pt}
\vfill


  \centerline{\Large The Six Isomorphism Classes of the Three-dimensional Algebras }

\vspace{.5cm}


\hspace{-2.4cm}$ \begin{array}{cccccccc} 
\mbox{\underline{algebra}} & \mbox{\underline{anti-rotor ($M_{\mathfrak{u}}$) }} & \mbox{\underline{unital norm group ($\mathfrak{u}^d$)}}  \\ 
\\
 \mathbb{R}\times \mathbb{R}\times \mathbb{R}  \approx \left[\begin{array}{ccc}
x  & 0 & 0   \\
 0 & y & 0 \\
 0 & 0 & z    
\end{array}\right]  \approx \mathbb{R}\times\slashed{\mathbb{C}}  &  \mathbb{R}\times M_{\mathfrak{u}_{\mathbb{R}\times\mathbb{R}}}, \begin{array}{ccc}
V_\mathbb{R}(\alpha\beta\gamma)   \\ 
(3,3,1,3,2,3)   
\end{array} & \begin{array}{ccc} x^{\frac{\alpha}{3}}y^{\frac{\beta}{3}}z^{\frac{\gamma}{3}} \\ (0,2,0) \end{array}
 \\ \\
   \mathbb{R}\times\mathbb{C} \approx \left[\begin{array}{ccc}
x  & 0 & 0   \\
 0 & y & -z \\
 0 & z & y    
\end{array}\right]   \hfill &  \mathbb{R}\times M_{\mathfrak{u}_{\mathbb{C}}}, \begin{array}{ccc}
V_\mathbb{R}(-\alpha(\beta^2+\gamma^2))   \\ 
(3,3,1,3,2,2)   
\end{array} & \begin{array}{ccc}  x^\alpha (y^2+z^2)^{\frac{\beta}{2}} \,e^{\gamma\,{\rm arctan}\left(\frac{z}{y}\right)} \\ (1,1,0)\end{array} \\ \\
   \mathbb{R}\times\mathbb{D} \approx \left[\begin{array}{ccc}
x  & 0 & 0   \\
 0 & y & z \\
 0 & 0 & y    
\end{array}\right]  \hfill &  \mathbb{R}\times M_{\mathfrak{u}_{\mathbb{D}}}, \begin{array}{ccc}
V_\mathbb{R}(-\alpha\gamma^2)   \\ 
(3,3,1,2,1,2)   
\end{array}  &  \begin{array}{ccc} x^\alpha y^{\beta} e^{\gamma\frac{y}{z}} \\ (0,1,1)\end{array} \\ \\
   \mathbb{T}_3  \approx \left[\begin{array}{ccc}
x  & y & z   \\
 0 & x & y \\
 0 & 0 & x    
\end{array}\right] \hfill &  \left[\begin{array}{ccc}
\alpha  & \beta & \gamma   \\
 \beta & \gamma & 0 \\
 \gamma & 0 & 0   
\end{array}\right], \begin{array}{ccc}
V_\mathbb{R}(\gamma^3)   \\ 
(3,3,1,1,0,1)   
\end{array}  & \begin{array}{ccc} x^\alpha e^{ \beta\frac{y}{x}  +\gamma \frac{z}{x} -\gamma \frac{ 1}{2}\left(\frac{y}{x}\right)^2 } \\ (0,0,2)\end{array} \\ \\
\left[\begin{array}{ccc}
x  & y & 0   \\
 0 & x & 0 \\
 0 & z & x    
\end{array}\right] \hfill &    \left[\begin{array}{ccc}
\alpha  & \beta & \gamma   \\
 \beta & 0 & 0 \\
 \gamma & 0 & 0   
\end{array}\right], \begin{array}{ccc}
V_\mathbb{R}(0)   \\ 
(3,2,1,0,3,0)   
\end{array}   &  \begin{array}{ccc} x^\alpha e^{ \beta\frac{y}{x}  +\gamma \frac{z}{x}  } \\ (0,0,2)\end{array} \\ \\
\left[\begin{array}{ccc}
x  & z & y   \\
 0 & x+y & 0 \\
 y & -z & x    
\end{array}\right]  \approx \left[\begin{array}{ccc}
x  & 0 & 0   \\
 0 & y & 0 \\
 0 & z & x    
\end{array}\right]  \hfill & \hspace{.5cm} \left[\begin{array}{ccc}
\alpha  & \beta & 0   \\
 \beta & \alpha & 0 \\
 0 & 0 & 0   
\end{array}\right] = M_{\mathfrak{u}_{\slashed{\mathbb{C}}}}\times \{0\}, \begin{array}{ccc}
V_\mathbb{R}(0)   \\ 
(2,2,1,0,2,0)
\end{array}  &  \begin{array}{ccc}(x+y)^\alpha (x-y)^\beta \\ (0,1,0)\end{array} \\ \\
\end{array} $



\vspace{.5cm}


\noindent {\bf Table II.} {\small Each row lists a representative  of one of the six three-dimensional algebra isomorphism classes along with its representation as a subalgebra of $M_3(\mathbb{R})$, followed by the associated anti-rotor and unital norm group.  $\mathbb{T}_3$ is the three-dimensional upper triangular Toeplitz matrix algebra. The invariant ordered sextuple in each row lists the value of the dimension of $\mathfrak{u}$, the largest rank of a member of $\mathfrak{u}$, the smallest rank of a nonzero member of $\mathfrak{u}$, the dimension of the parameter subspace to which det$(M_{\mathfrak{u}})$ is sensitive, the real dimension of the real algebraic set defined by det$(M_{\mathfrak{u}}) =0$, and the number of irreducible components of $V_\mathbb{R}(\text{det}(M_{\mathfrak{u}}))$ (this hypersurface is indicated above the sextuple).  The reduced triple is indicated below each unital norm group.}
 
\vfill
\hspace{0pt}
\pagebreak

\section{More examples, and application to some not associative algebras}
\subsection{Upper triangular Toeplitz matrix algebras}

The anti-rotor and unital norm group associated with an upper triangular Toeplitz matrix algebra have some interesting features.
\begin{itemize}\item The anti-rotor is given by the left regular representation (the algebra whose elements are lower triangular Toeplitz matrices) multiplied on the left by the exchange matrix. \item The usual norm is sensitive to only a single component of an element, but almost all unital norms are sensitive to all components of an element.\end{itemize}


 \begin{definition}\label{5/4/22.1} $\mathbb{T}_n$ is the $n$-dimensional vector space of real upper triangular Toeplitz $(n\times n)$-matrices. For $s\in \mathbb{T}_n$, the notation $s=(x_1,\dots,x_n)$ indicates that $s$ corresponds to the Toeplitz upper triangular $(n\times n)$-matrix whose entries on the successive diagonals, beginning with the main diagonal and proceeding to the right, are given by the successive vector component values $x_1,\dots,x_n$.   
\end{definition}

In the sequel, for $s\in \mathbb{T}_n$, $s$ can refer either to its format as an $n$-vector or its format as an $(n\times n)$-matrix, depending on the context.


The following result is known, but we include the proof for completeness, and to introduce some notation.  



 \begin{proposition}\label{5/10/22.1} Under the usual matrix product, $\mathbb{T}_n$ is a real unital commutative associative $n$-dimensiona algebra.\end{proposition} 
 
 \begin{proof}
Let ${\bf z}$ be the $(n \times n)$-matrix with entries having value $1$ on its superdiagonal  (the diagonal immediately to the right of the main diagonal) and value $0$ otherwise.  It is easily shown that for $j<n$ its $j$-th power, ${\bf z}^j$, is the $(n\times n)$-matrix with entries $1$ on the $j$-th diagonal to the right of the main diagonal, with all other entries being $0$.   Thus, we can express the matrix format of $s\in\mathbb{T}_n$ as a polynomial in ${\bf z}$, such that
\begin{equation}\label{5/11/22.5} s = x_1 {\bf 1} + x_2 {\bf z} + x_3 {\bf z}^2 + \cdots + x_n {\bf z}^{n-1}.\end{equation}  It is similarly easy to show that ${\bf z}^k$ is the matrix of all zero entries if $k\ge n$.  It follows that for $s,t\in \mathbb{T}_n$, we have $st\in \mathbb{T}_n$.   Based on the polynomial format of an element as given in (\ref{5/11/22.5}), the proposition follows.
 \end{proof}
 

We will show that upper triangular Toeplitz matrix algebras have unital norms sensitive to all element components, providing a great contrast with the usual norm for these algebras, which is sensitive to only one element component.  One can easily generate other algebras with unital norms ``similarly dominating" the usual norm, as seen later.

To approach the problem of deriving the anti-rotor for $\mathbb{T}_n$, the following will be important.

\begin{lemma} \label{5/7/22.15} For a unit $s\in \mathbb{T}_n$,  $n>1$, let $\left( s^{-1} \right)_p$ be the $p$-th component of $s^{-1}$ in the latter's $n$-vector format.  Given $k\le n$, then for all $j<k$ we have \begin{equation}\label{5/6/22.10} \frac{\partial \left( s^{-1} \right)_k}{\partial x_j} = \frac{\partial \left( s^{-1} \right)_{k-j+1}}{\partial x_1}.\end{equation}
\end{lemma}

\begin{proof} The proof is by induction. Equation (\ref{5/6/22.10}) is a tautology for the $k=2$ case.  We thus assume the inductive hypothesis that (\ref{5/6/22.10}) is true  for all cases up through $k-1$, and then establish its truth for the $k$-th case.


We must first express $s^{-1}$ in terms of the components of  $s=(x_1,\dots,x_n)\in \mathbb{T}_n$ with $x_1\ne 0$.   From the proof of Proposition \ref{5/10/22.1},  $s^{-1}$ is an upper triangular Toeplitz matrix.  But as an $n$-vector, it is easily seen that we can write it as  \begin{equation}\label{5/9/22.1} s^{-1} = (\frac{1}{x_1}, \frac{y_1}{x_1},\dots, \frac{y_{n-1}}{x_1}),\end{equation} where the $y_i$ are to be determined. 
Based on the notation introduced in the proof of Proposition \ref{5/10/22.1}, we can use (\ref{5/11/22.5}) to write 
\begin{align}\nonumber s &= x_1 \left({\bf 1} + \frac{x_2}{x_1} {\bf z} + \frac{x_3}{x_1} {\bf z}^2 + \cdots + \frac{x_n}{x_1} {\bf z}^{n-1}\right)\\ s^{-1} &= x_1^{-1} \left({\bf 1} + y_1 {\bf z} + y_2 {\bf z}^2 + \cdots + y_{n-1} {\bf z}^{n-1}\right).  \label{5/18/22.10} \end{align}  Substituting the above two equations into $s^{-1} s = {\bf 1}$ leads to following set of relations, which allow the $y_i$ to be computed recursively:
\begin{align} -y_1 &= \frac{x_2}{x_1} \nonumber \\ -y_2 &= \frac{x_2}{x_1}y_1 + \frac{x_3}{x_1} \nonumber \\ &\;\vdots \nonumber \\ -y_{k-1} &= \frac{x_2}{x_1} y_{k-2} + \frac{x_3}{x_1} y_{k-3} + \cdots + \frac{x_k}{x_1} \label{5/7/22.1} \\ &\;\vdots \nonumber \end{align} 
 
 \noindent For $k>1$, (\ref{5/18/22.10}) implies  $\left( s^{-1} \right)_k = \frac{y_{k-1}}{x_1}$.  Applying this to (\ref{5/7/22.1}) with $k>1$, we find
\begin{eqnarray} \nonumber \left( s^{-1} \right)_k &=& -\frac{x_2}{x_1}\left( s^{-1} \right)_{k-1} - \frac{x_3}{x_1}\left( s^{-1} \right)_{k-2} -\dots  - \frac{x_k}{x_1}\left( s^{-1} \right)_1 \\ &=& -\frac{1}{x_1}\sum_{m=1}^{k-1} x_{m+1}\left( s^{-1} \right)_{k-m}. \label{5/7/22.2}\end{eqnarray}
Differentiating with respect to $x_j$, we find \begin{equation}  \frac{\partial \left( s^{-1} \right)_k}{\partial x_j} = -\frac{1}{x_1} \sum_{m=1}^{k-j} x_{m+1} \frac{\partial \left( s^{-1} \right)_{k-m}}{\partial x_j} -  \frac{1}{x_1}\left( s^{-1} \right)_{k-j+1} \label{5/7/22.4}, \end{equation} where the upper limit on the sum on the right-hand-side is $(k-j)$ rather than $(k-1)$ because the partial derivatives in the sum are zero for $m> k-j$ (i.e., $\left( s^{-1} \right)_p$ is not a function of $x_j$ for $p<j$).  

Next, substituting $(k-j+1)$ for $k$ in (\ref{5/7/22.2}) and then differentiating with respect to $x_1$, we have   \begin{eqnarray} \nonumber
 \frac{\partial \left( s^{-1} \right)_{k-j+1}}{\partial x_1} &=& \frac{1}{x_1^2} \sum_{m=1}^{k-j} x_{m+1} \left( s^{-1} \right)_{k-j-m+1}    - \frac{1}{x_1} \sum_{m=1}^{k-j} x_{m+1} \frac{\partial \left( s^{-1} \right)_{k-j-m+1}}{\partial x_1}  \\ &=& - \frac{1}{x_1}\left( s^{-1} \right)_{k-j+1}  -  \frac{1}{x_1} \sum_{m=1}^{k-j} x_{m+1} \frac{\partial \left( s^{-1} \right)_{k-j-m+1}}{\partial x_1},\label{5/7/22.3}\end{eqnarray} where the equality of the first terms on the right-hand-sides of the two above equalities comes from (\ref{5/7/22.2}) itself (with $(k-j+1)$ substituting for $k$). 
 Now, according to the induction hypothesis, $\frac{\partial \left( s^{-1} \right)_{k-m}}{\partial x_j} = \frac{\partial \left( s^{-1} \right)_{k-j-m+1}}{\partial x_1}$.  Inserting this into (\ref{5/7/22.3}), we see that the left-hand-sides of (\ref{5/7/22.4}) and (\ref{5/7/22.3}) are equal, as was required to be shown.
 \end{proof} 

Let $\slashed{\bf 1}_n$ be the $(n\times n)$ exchange matrix (the matrix whose entries are 1 on the antidiagonal and zero everywhere else).

\begin{theorem} \label{5/4/22.2} For any $t\in \mathbb{T}_n$ let $\hat{t}$ be its $(n\times n)$-matrix format.    Then $\hat{t} \slashed{\bf 1}_n$ is an uncurling metric.  Conversely, each member of the anti-rotor associated with $\mathbb{T}_n$ can be expressed as $\hat{t}\slashed{\bf 1}_n$ for some $t\in \mathbb{T}_n$.   \end{theorem}


\begin{proof} The following notation is more convenient for the moment.  Given $(\gamma_n,\dots,\gamma_1)$ $\in \mathbb{T}_n$, let $G_n$ be its upper triangular Toeplitz $(n\times n)$-matrix format (which happens to be the transpose of its left regular representation), and consider the resulting Hankel matrix, 
\begin{equation}\label{5/7/22.20} H_n \equiv G_n \slashed{\bf 1}_n = \begin{bmatrix} \gamma_1 & \gamma_2 & \gamma_3 & \dots & \gamma_n \\ \gamma_2 & \gamma_3 & \dots & \gamma_n & 0 \\ \gamma_3 & \dots & \iddots & \iddots & 0  \\ \vdots & \iddots & \iddots & \iddots & \vdots \\ \gamma_n & 0 & \dots & \ldots & 0 \end{bmatrix}.\end{equation}
For $\left( s^{-1} \right)_p$ as defined in the statement of Lemma \ref{5/7/22.15}, we then have
\begin{multline}  H_n s^{-1}   =  \Big( \gamma_1\left( s^{-1} \right)_1  + \dots + \gamma_n \left( s^{-1} \right)_n, \\ \gamma_2\left( s^{-1} \right)_1  + \dots + \gamma_n \left( s^{-1} \right)_{n-1}, 
  \dots , 
 \gamma_n \left( s^{-1} \right)_1 \Big). \label{5/7/22.40}
\end{multline}
 We will first demonstrate that $H_n$ is an uncurling metric.  According to (\ref{4/7/23.1}) as referred to by Definition \ref{9/20/23.1}, we are required to demonstrate 
\begin{equation}\label{5/7/22.30} d\left([H_n s^{-1}]\cdot \mathfrak{d}s\right) = 0.\end{equation}  

To verify (\ref{5/7/22.30}), we apply an induction argument.  The $n=2$ case is easily demonstrated to satisfy the above equation (e.g., see Section 4.4).  We assume the induction hypothesis that (\ref{5/7/22.30}) holds for all cases less than $n$, and now need to show that it holds for the $n$-th case.  
We begin by showing that the partial derivative of the first component of $H_n s^{-1}$ with respect to any $x_j$, with $j>1$, minus the partial derivative of the $j$-th component of $H_n s^{-1}$ with respect to $x_1$, is zero. We will then show that the analogous difference of partial derivatives of {\it any} two different components of $H_ns^{-1}$ is also zero.

So, from (\ref{5/7/22.40}), 
\begin{equation}\label{5/18/22.5} \frac{\partial \left(H_n s^{-1}\right)_j}{\partial x_1} = 
\sum_{i=j}^n \gamma_i\frac{\partial \left(s^{-1}\right)_{i-j+1}}{\partial x_1}. 
\end{equation} \enlargethispage{1.5cm}
But, \begin{equation}\label{5/18/22.6}  \frac{\partial \left(H_n s^{-1}\right)_1}{\partial x_j} = \sum_{i=1}^n \gamma_i\frac{\partial \left(s^{-1}\right)_i}{\partial x_j}
= \sum_{i=j}^n \gamma_i\frac{\partial \left(s^{-1}\right)_i}{\partial x_j} = \sum_{i=j}^n \gamma_i\frac{\partial \left(s^{-1}\right)_{i-j+1}}{\partial x_1},
\end{equation} where the second equality follows (\ref{5/9/22.1}) and (\ref{5/7/22.1}), which indicate that $\frac{\partial \left(s^{-1}\right)_i}{\partial x_j} = 0$ for $i<j$, and the third equality follows from Lemma \ref{5/7/22.15}.
Equations (\ref{5/18/22.5}) and (\ref{5/18/22.6})  demonstrate that the partial derivative of the first component of $H_n s^{-1}$ with respect to any $x_j$, with $j>1$, minus the partial derivative of the $j$-th component of $H_n s^{-1}$ with respect to $x_1$, is equal to zero.

Now we will consider the vanishing of the partial derivative of the $k$-th component of $H_n s^{-1}$ with respect to $x_j$ minus the partial derivative of the $j$-th component of $H_n s^{-1}$ with respect to $x_k$, for all possible $k>j>1$ (i.e, all choices relevant to (\ref{5/7/22.30}) other than considered in the prior paragraph).  In fact, based on the telescoping format of (\ref{5/7/22.40}), we already have the result that this difference vanishes, because it merely reduces to an instance of the theorem for the case of a matrix of the form $H_ks^{-1}$ with $k<n$ - and this case is assumed to hold according to the induction hypothesis.  So, (\ref{5/7/22.30}) is established.

Furthermore, a Hankel matrix of the form (\ref{5/7/22.20}) is the only possibility for an uncurling metric $L$ because any other choice resulting in the coefficients of the terms in any two components of $Ls^{-1}$ cannot lead to a null value for the difference of the two requisite partial derivatives (as is necessary for $d\left([L s^{-1}]\cdot \mathfrak{d}s\right) = 0$) due to linear independence of the individual terms summing to the components. That is, replacing  $H_n$ with matrix $L=(a_{p,q})_{p,q=1}^n$ in (\ref{5/18/22.5}) 
and (\ref{5/18/22.6}), and invoking Lemma \ref{5/7/22.15}, we obtain
\begin{eqnarray}\nonumber \frac{\partial \left(L s^{-1}\right)_j}{\partial x_1}  -  \frac{\partial \left(L s^{-1}\right)_1}{\partial x_j} &=&  \sum_{i=1}^n a_{j,i}\frac{\partial \left(s^{-1}\right)_i}{\partial x_1} - \sum_{i=j}^n a_{1,i}\frac{\partial \left(s^{-1}\right)_i}{\partial x_j} \\  &=& \sum_{i=1}^n a_{j,i}\frac{\partial \left(s^{-1}\right)_i}{\partial x_1} - \sum_{i=j}^n a_{1,i}\frac{\partial \left(s^{-1}\right)_{i-j+1}}{\partial x_1} \nonumber \\
&=& \sum_{i=1}^n a_{j,i}\frac{\partial \left(s^{-1}\right)_i}{\partial x_1} - \sum_{i=1}^{n-j+1} a_{1,i+j-1}\frac{\partial \left(s^{-1}\right)_i}{\partial x_1}.
\label{5/9/22.10}
\end{eqnarray}
This indicates that we must have $a_{j,i}=0$ for $i>n-j+1$, and otherwise have $a_{j,i} = a_{1,i+j-1}$.  Since this must hold for any $j>1$, $L$ must be a Hankel matrix of the form (\ref{5/7/22.20}).  Thus, the $H_n$ are the only possibilities for uncurling metrics. 

Since $t=(\gamma_n,\dots,\gamma_1)$ is an arbitrary member of $\mathbb{T}_n$, the theorem follows.
\end{proof}

\begin{corollary}\label{7/1/22.1} For any $t\in \mathbb{T}_n$, denote the value of the $(1,n)$ entry of its matrix format $\hat{t}$ as $\hat{t}_{1,n}$.  Then $\frac{1}{\hat{t}_{1,n}}\hat{t}\slashed{\bf 1}_n$ is a normalized uncurling metric if $\hat{t}_{1,n}\ne 0$, and all normalized uncurling metrics can be written in this form for some $t\in \mathbb{T}_n$. \end{corollary}
\begin{proof} Theorem \ref{5/4/22.2} has shown that all uncurling metrics of $\mathbb{T}_n$ have the form (\ref{5/7/22.20}).
We now consider constraint (\ref{7/7/21.1}) of Definition \ref{9/20/23.1}.  Let $H_n^{(1)}$ be matrix $H_n$ except with the  $(1,1)$-entry having the value $1$.  \enlargethispage{2cm}
Since $\|{\bf 1}\|^2 =1$ for this algebra, we must show that \[ s^T H_n^{(1)} s^{-1} = s\cdot [H_n^{(1)}s^{-1}] = 1.\] From (\ref{5/9/22.1}) and (\ref{5/7/22.20}), we have
\begin{eqnarray}\nonumber s^T H_n s^{-1} &=& (x_1,\dots,x_n)\cdot \left[ H_n \left(\frac{1}{x_1}, \frac{y_1}{x_1},\dots,\frac{y_{n-1}}{x_1}\right)^T\right] \\ 
 &=&  \gamma_1 + \gamma_2y_1 + \gamma_3y_2 \hspace{3mm} + \hspace{3mm}\gamma_4y_3 \hspace{3mm} + \dots   \hspace{6mm}+\gamma_ny_{n-1} \nonumber \\
&+&  \gamma_2 \frac{x_2}{x_1} + \gamma_3\frac{x_2}{x_1}y_1 + \hspace{2mm} \gamma_4\frac{x_2}{x_1}y_2 + \dots +  \gamma_n\frac{x_2}{x_1}y_{n-2} \nonumber \\ &+&\gamma_3 \frac{x_3}{x_1}  + \hspace{3mm}\gamma_4\frac{x_3}{x_1}y_1 +  \dots \hspace{4mm} +  \gamma_n\frac{x_3}{x_1}y_{n-3} \nonumber \\ 
&+&  \hspace{5mm} \vdots \hspace{2.5cm} \iddots \bigskip\bigskip \nonumber \\ &+& \gamma_{n-1} \frac{x_2}{x_1} \hspace{3mm} + \gamma_n \frac{x_2}{x_1}y_1 \nonumber \\ &+& \hspace{2mm} \gamma_n\frac{x_n}{x_1}. \label{6/23/22.1}
 \end{eqnarray}
 For any $i>1$, if we collect all the terms on the right-hand-side above that each have $\gamma_i$ as a coefficient, and then apply (\ref{5/7/22.1}), it is quickly seen that the sum of all terms having $\gamma_i$ as a coefficient vanishes.  It follows that
 $s^T H_n s^{-1} =\gamma_1$.   Thus, the requisite matrix satisfying constraint (\ref{7/7/21.1}) of Definition \ref{9/20/23.1} will be the right-hand-side of (\ref{5/7/22.20}) with $\gamma_1=1$.  Setting $\frac{1}{\hat{t}_{1,n}}t = (\gamma_n,\dots,\gamma_2, 1)$, the corollary follows.
\end{proof}

\begin{corollary}\label{5/17/22.1} For $s= (x_1,\dots,x_n)\in \mathbb{T}_n$ with $x_1\ne 0$, the usual norm is $x_1^n$.  The special unital norm family of $\mathbb{T}_n$ is the $(n-1)$-parameter family,
\begin{equation} \label{5/5/22.6} \ell_{\gamma_2,\dots,\gamma_n}(x_1,\dots,x_n) = x_1 e^{P_{\gamma_2,\dots,\gamma_n}\left(\frac{x_2}{x_1},\dots,\frac{x_n}{x_1}\right)},\end{equation} where $P_{\gamma_2,\dots,\gamma_n}(w_1,\dots,w_{n-1})$ is a real $(n-1)$-degree multivariate polynomial with coefficients dependent on the $(n-1)$ real parameters $\gamma_2,\dots,\gamma_n$.    \end{corollary}

\begin{proof} The usual norm of $s=(x_1,\dots,x_n)\in \mathbb{T}_n$ is immediately seen to be $x_1^n$.

With respect to the notation $G_n,H_n$ introduced at the beginning of the proof of Theorem \ref{5/4/22.2}, each $(\gamma_n,\dots,\gamma_2,1)\in\mathbb{T}_n$ is associated with a Toeplitz matrix $G_n^{(1)}$ having its $(1,n)$-entry equal to $1$.  Corollary \ref{7/1/22.1} indicates that $H^{(1)}_n\equiv G_n^{(1)}\slashed{\bf 1}_n$ is a normalized uncurling metric. From (\ref{11/7/20.2}),  the associated unital norm is $e^{\left[\int_{\bf 1}^s (H_n^{(1)} v^{-1})\cdot d\mathfrak{v}\right]}$.  Elementary integration and a simple induction argument quickly results in (\ref{5/5/22.6}) and the $(n-1)$-degree of the multivariate polynomial in question.
\end{proof}
 
 \subsection{Extension to certain not associative algebras}
The two-dimensional real unital algebras (there are three, up to an isomorphism) have generalizations as the Cayley-Dixon sequence algebras ($\mathbb{R}$ and $\mathbb{C}$ being its lowest dimensional members),  the Spin Factor Jordan algebras (generalizing $\slashed{\mathbb{C}}$), and real upper triangular matrix algebras with at least some redundant diagonal components (generalizing the lowest dimensional case, $\mathbb{D}$).  Their unital norm groups will be the subject of this subsection. 

Spin Factor Jordan Algebras and Cayley-Dixon Algebras after the quaternions are power-associative but not associative.  While the usual norm is thereby inapplicable, there is an ``Algebraic Norm" replacement \cite{jacobson:1963} (based on the minimal polynomial, and thus analogous to the usual norm as given by the constant term of a characteristic polynomial).  For the Spin Factor Jordan Algebras and Cayley-Dixon Algebras, the Algebraic Norms are essentially the  Minkowski  norm and Euclidean norm, respectively.  Thus, useful norms are not restricted to associative (or alternative) algebras, even though it turns out that the multiplicative property (i.e., $N(s_1s_2)=N(s_1)N(s_2)$) must then be sacrificed  \cite{schafer:1963}. 
So, despite a strong bias against a norm not satisfying the multiplicative property, these examples indicate that ``algebraic" norms that are not multiplicative can nevertheless be very significant \cite{mccrimmon:2004}. 

Of course, apart from the omnipresent unital norm group member comparable to the usual norm as guaranteed by Theorem \ref{11/17/22.1} for real finite-dimensional unital associative algebras (as well as all unital norms resulting from two-dimensional real unital algebras), unital norms are not typically multiplicative.  The multiplicative property is instead replaced by (\ref{5/23/23.3}).  The latter equation also replaces the multiplicative property for the Algebraic Norm in the cases of the not associative Spin Factor Jordan Algebras and Cayley-Dixon Algebras. 
Like the Algebraic Norm, the unital norm/uncurling metric constructions remain valid for these latter two classes of algebras.  The way we deal with these particular not associative algebras is as follows.

\subsubsection{Unital norms for $^*$-algebras whose Hermitian elements are real} 

Consider an algebra associated with an involutive antiautomorphism ``
$^*$ " (known as a $^*$-algebra, star-algebra, or involutive algebra) such that the Hermitian elements are members of $\mathbb{R}$ (thus, $s+s^*$ and $s^*s$ are both real; also, $s^*s = ss^*$ because $(s+s^*)s = s(s+s^*)$ since $(s+s^*)$ is real). 
Then \[ X^2 - (s+s^*)X + s^*s, \] is the minimal polynomial of any not real element $s$, and annihilates all elements, real and not real.  Consistent with comments in the second paragraph of Section 4.2, we define its constant term to be the ``Algebraic Norm" of an element $s$, by analogy with the characteristic polynomial defined for an associative algebra, which supplies the norm from the constant term of that element-associated polynomial annihilating all elements.


On these $^*$-algebras, the multiplicative identity $\mathbf{1}$ is $1$ (where the latter is thought of as the algebra element that is the multiplicative identity on the subalgebra $\mathbb{R}$ rather than a field element).  Accordingly, for these algebras, we define the expression $\|{\bf 1}\|^2$ (e.g., that appears in Definition \ref{9/20/23.1}) to be $1$.  We also have,

\begin{lemma}\label{9/14/22.2} The inverse of a unit $s$ in a finite-dimensional $^*$-algebra whose Hermitian elements are real is unique, and given by $s^{-1}=\frac{s^*}{ss^*}$.\end{lemma}

\begin{proof} 
We can write $s=\frac{s+s^*}{2} +\frac{s-s^*}{2}$, where the first  term on the right-hand-side is real and the second term on the right-hand-side is antisymmetric with respect to application of the $^*$-operation.  We define $\sigma$ to be the first term on the right-hand-side, and $\bar{s}$ to be the second term on the right-hand-side.  Suppose $x = \chi + \bar{x}$ is an inverse of $s$, where $\chi$ is real, and $\bar{x}$ is antisymmetric. Then \begin{equation}\label{9/14/22.1} sx = \sigma\chi + \sigma\bar{x}+\chi\bar{s} + \bar{s}\bar{x} = 1,\end{equation} and $xs = \sigma\chi + \sigma\bar{x}+\chi\bar{s} + \bar{x}\bar{s} = 1$.  Thus, $\bar{s}\bar{x} = \bar{x}\bar{s}$.  Since $\bar{s}^* = (\frac{s-s^*}{2})^* = -(\frac{s-s^*}{2}) = -\bar{s}$, we have $(\bar{s}\bar{x})^* = \bar{x}^*\bar{s}^* = \bar{x}\bar{s} = \bar{s}\bar{x}$.  Hence, $\bar{s}\bar{x}$ is real.  It follows from (\ref{9/14/22.1}) that $\sigma\chi+\bar{s}\bar{x} =1$ and $\sigma\bar{x}+\chi\bar{s} = 0$.  The latter has has a unique solution, yielding $\chi+\bar{x} = \frac{\sigma-\bar{s}}{\sigma^2-\bar{s}\bar{s}} = \frac{s^*}{ss^*}$, where the second equality follows since the right-hand-side is obviously an inverse ($1=s\frac{s^*}{ss^*} = \frac{s^*}{ss^*}s$ because $s^*s=ss^*$ is real and thus self-adjoint), and the inverse has just now been seen to be unique.
\end{proof}

Extending Definition \ref{9/20/23.1} to include this class of $^*$-algebras, we can now provide a special unital norm  by identifying a linear transformation $L$ satisfying (\ref{4/7/23.1}), (\ref{7/7/21.1}) and then evaluating (\ref{11/7/20.2}).  
 
 \begin{proposition}\label{1/11/22.1} Given a $^*$-algebra with vector space of elements $\mathbb{R}^n$ and whose Hermitian elements are members of $\mathbb{R}$, the Algebraic Norm on the $^*$-algebra in a sufficiently small neighborhood of ${\bf 1}$ is the square of a special unital norm. \end{proposition}

\begin{proof} By definition, the Algebraic Norm is \begin{equation}\label{1/15/22.1} N(s)= ss^* = N(s^*)\in\mathbb{R}.\end{equation} 
 We then have $N(\alpha s) = \alpha^2 N(s)$ for $\alpha\in \mathbb{R}$, i.e., degree-2 positive homogeneity of $N(s)$.  
For $N(s)\ne 0$, Lemma \ref{9/14/22.2} implies 
\begin{equation}\label{1/12/22.1} s^{-1} = \frac{s^*}{N(s)}.\end{equation}
Applying $N(\cdot)$ to both sides of  (\ref{1/12/22.1}), the degree-2 positive homogeneity and (\ref{1/15/22.1}) imply $N(s)N(s^{-1})=1$, so that \begin{equation}\label{1/15/22.2} \sqrt{N(s)}\sqrt{N(s^{-1})}=1. \end{equation}

We are required to show that $\sqrt{N(s)}$ is a special unital norm.  Thus, we must first verify that $\sqrt{N(s)}$ satisfies the requirements of the extended  version of Definition \ref{9/20/23.1} as indicated prior to the statement of this proposition. 

To begin with, (\ref{1/15/22.2}) already satisfies  (\ref{5/23/23.3}).
 Equation (\ref{1/15/22.1}) indicates that $N(s)$ is a quadratic form, so there is an orthonormal basis with respect to which the members of the algebra can be written as $s= (s_1,\dots, s_i,\dots)$ with $N(s) = \sum_i \delta_i s_i^2$ and $\delta_i = -1$, $0$, or $1$, since $^*$ is an involution.  Let $K_1 \equiv {\rm diag}\{\delta_i\}$.   Then, \begin{equation}\label{2/6/22.1} N(s) = s\cdot(K_1 s).\end{equation}
If $s^{-1}$ exists (i.e., $N(s)\ne 0$), we then have \begin{equation}\label{2/6/22.2} \nabla\left[\sqrt{N(s)}\right] = \frac{K_1s}{\sqrt{N(s)}}.\end{equation}  Define $K_2$ to be the algebra's involutive antiautomorphism operation $^*$.  It follows that $K_1K_2K_2[\frac{s}{\sqrt{N(s)}}] = \nabla\left[\sqrt{N(s)}\right]$. Defining $L\equiv K_1K_2$, dividing both sides of the equation in the last sentence by $\sqrt{N(s)}$, and using (\ref{1/15/22.1}) and (\ref{1/12/22.1}), we then have \begin{equation}\label{1/12/22.2} Ls^{-1}  = \frac{\nabla\left[\sqrt{N(s)}\right]}{\sqrt{N(s)}} .\end{equation} This equation  is valid for $N(s)>0$, and the latter pertains in a sufficiently small neighborhood of ${\bf 1}$.  Thus,  (\ref{11/7/20.1}) is satisfied.  From (\ref{2/6/22.1}), (\ref{2/6/22.2}), and (\ref{1/12/22.2}), we have \[ s\cdot [Ls^{-1}] = s\cdot K_1\left[\frac{s}{N(s)}\right] = 1,\] which is (\ref{7/7/21.1}).  
 
 To establish that $\sqrt{N(s)}$ is a unital norm, we  now only need to verify that (\ref{4/7/23.1}) is satisfied, i.e., $d\left([Ls^{-1}]\cdot \mathfrak{d}s\right)=0$.  According to (\ref{2/6/22.2}), (\ref{1/12/22.2}), and the definition  of $K_1$, we have $d\left([Ls^{-1}]\cdot \mathfrak{d}s\right) = d\left(\sum_i \frac{\delta_is_i ds_i}{N(s)}\right)$.  Performing the requisite differentiation to obtain the component of the exterior derivative associated with a choice of $i$ and $j$, we find \[ \frac{\partial [Ls^{-1}]_j}{\partial s_i} - \frac{\partial [Ls^{-1}]_i}{\partial s_j} = -2\frac{(\delta_j s_j)(\delta_i s_i)}{N^2(s)} +
2\frac{(\delta_i s_i)(\delta_j s_j)}{N^2(s)} = 0,  \] where $\frac{\partial [Ls^{-1}]_j}{\partial s_i}$ is the partial derivative of the $j$-th component of $[Ls^{-1}]$ with respect to the $i$-th component of $s$, etc. Thus,  $d\left([Ls^{-1}]\cdot \mathfrak{d}s\right) = 0$.

Consequently, $\sqrt{N(s)}$ satisfies the special unital norm requirements of the extended version of Definition \ref{9/20/23.1}.  
\end{proof}


\subsubsection{The Cayley-Dixon sequence}  From the proof of Proposition \ref{1/11/22.1}, it is easily shown that a normalized uncurling metric associated with an algebra of the Cayley-Dixon sequence (all of whose members are $^*$-algebras whose Hermitian elements are real) is given by ${\rm diag}\{1,-1,-1,\dots,-1\}$, with the associated special unital norm seen to be the Euclidean norm. Although it is easily shown from results in \mbox{Section 9} that $\mathbb{C}$ is associated with a 1-parameter family of normalized uncurling metrics, a simple argument (using Lemma \ref{9/14/22.2}) indicates that all other algebras of the sequence have only the above normalized uncurling metric, i.e., the special unital norm family has only one member.  
 
\subsubsection{Spin Factor Jordan Algebras}  Consider the bilinear product ``$\,\bullet\,$" such that for
$\alpha,\beta\in \mathbb{R},\,\, {\sf a},{\sf b}\in \mathbb{R}^m$,  \begin{equation}  (\alpha+ {\sf a})\bullet (\beta + {\sf b}) \equiv
    [\alpha \beta + {\sf a} \cdot {\sf b}] + [ \beta {\sf a} + \alpha {\sf b}] \in \mathbb{R}\oplus \mathbb{R}^m. \label{12/24/18.1}\end{equation}  This makes the vector space $\mathbb{R}\oplus \mathbb{R}^m$ a unital algebra with the multiplicative identity ${\bf 1}$ as simply $1$ (where in this instance $1$ is understood to be an element of the algebra rather than a member of the field).  For $m\ge 2$, this is known as a Spin Factor Jordan Algebra \cite{mccrimmon:2004}.

 There are two involutive antiautomorphisms on each of these commutative algebras, one of which is the identity and the other is the conjugate operation, $(\alpha+ {\sf a})^* = (\alpha- {\sf a})$. The latter is the only one of the two involutive antiautomorphisms  for which the Hermitian elements are real.  Applying Proposition \ref{1/11/22.1}, it is easily shown that the special unital norm is the Minkowski norm.  Similar to the case of the Cayley-Dixon algebras, it is easy to show that this is the only special unital norm.

 
 \subsection{The algebra with vector space of elements $\mathbb{R}^n$ and component-wise multiplication} A higher dimensional analogue of the direct product of rings \mbox{$\mathbb{R}\times \mathbb{R}$}, i.e., $\prod_{i=1}^n \mathbb{R}$, is associated with an $(n-1)$-parameter family of  special unital norms.

 This algebra is isomorphic to the subalgebra of $ M_n(\mathbb{R})$  whose members are diagonal matrices. In this example, ${\bf 1} = (1,\dots,1)$ and $\|{\bf 1}\|^2 = n$.   It is easily appreciated that (\ref{4/7/23.1}) forwards the anti-rotor ${\rm diag} \{ \sigma_1, \sigma_2,\dots,\sigma_n\}$, so the anti-rotor is the same as the members of the left regular representation of the algebra.  This means that for $s=(x_1,\dots,x_n)$, we must have $L[(x_1,x_2,\dots,x_n)^{-1}] = (\frac{\sigma_1}{x_1}, \frac{\sigma_2}{x_2},\dots,\frac{\sigma_n}{x_n})$.     Evaluation of (\ref{11/7/20.2}) indicates that the members of the unital norm group are given by  \begin{equation}\label{1/15/22.5} \ell_{\sigma_1,\dots,\sigma_n}(x_1,\dots,x_n) = (x_1^{\sigma_1} x_2^{\sigma_2} \dots x_n^{\sigma_n})^\frac{1}{n}. \end{equation} From (\ref{7/7/21.1}), the special unital norms are given by (\ref{1/15/22.5}) under the constraint $\sum_1^n \sigma_i = n$. Thus, for this algebra we have a $(n-1)$-dimensional special unital norm family.   
 
 The geometric mean is the special unital norm given by (\ref{1/15/22.5}) with $\sigma_1=\sigma_2=\dots =\sigma_n = 1$.  
   
 \subsection{Upper triangular matrix algebras with independent entries on the main diagonal} Consider the real three-dimensional algebra with elements $(x,y,z)$ which multiply according to the matrices 
$\left[
\begin{array}{ccc}
 x & z   \\
 0 & y    
\end{array}
\right]$.  Or in other words, the algebra is $\left\{ \left[
\begin{array}{ccc}
 x & z   \\
 0 & y    
\end{array}
\right],x,y,z\in\mathbb{R}\right\}$.  We have $(x,y,z)^{-1} = \left( \frac{1}{x}, \frac{1}{y}, -\frac{z}{xy} \right)$, and ${\bf 1}= (1,1,0)$, and $\|{\bf 1}\|^2 = 2$.  Starting with a generic symmetric $L = \left[
\begin{array}{ccc}
 \alpha & \beta & \gamma   \\
 \beta & \delta & \epsilon \\ \gamma & \epsilon & \xi    
\end{array}
\right]$, 
it is required that the curl of $L(x,y,z)^{-1}$, which is \[\left( \frac{\epsilon}{xy} - \frac{\epsilon}{y^2} + \frac{\xi z}{xy^2},\,\,\, \frac{\gamma}{x^2} - \frac{\gamma}{xy} - \frac{\xi z}{x^2y}, \,\,\, \frac{\beta}{y^2} - \frac{\beta}{x^2} - \frac{\gamma z}{xy^2} + \frac{\epsilon z}{x^2y} \right),\] be such that each of its components vanishes in a neighborhood of $(1,1,0)$.  This implies that all entries of the matrix $L$ must be zero except for $\alpha$ and $\delta$.   Thus, $L(x,y,z)^{-1} = (\frac{\alpha}{x}, \frac{\delta}{y},0)$.  From (\ref{7/7/21.1}), we have the requirement that a normalized uncurling metric must satisfy \[ (x,y,z)\cdot [L(x,y,z)^{-1}] = \alpha +\delta =  2.\]  Equivalently, we can alternatively specify $\alpha \equiv 1+\sigma$ and $\delta \equiv 1-\sigma$, for arbitrary $\sigma\in\mathbb{R}$.   We thus have the normalized uncurling metric family composed of all  \begin{equation}\label{11/17/22.2} L = \left[
\begin{array}{ccc}
 1+\sigma & 0 & 0   \\
 0 & 1-\sigma & 0 \\ 0 & 0 & 0    
\end{array}
\right],\end{equation} $\sigma\in\mathbb{R}$.  In this case, each member of the family is singular.  Equation (\ref{11/7/20.2}) then leads to the 1-parameter family of special unital norms, \begin{equation}\label{12/21/21.1} \ell_\sigma(x,y,z) = x^\frac{1+\sigma}{2} y^\frac{1-\sigma}{2} = (xy)^{\frac{1}{2}}\left(\frac{x}{y}\right)^{\frac{\sigma}{2}}.\end{equation} 

The left regular representation of an element $(x,y,z)$ is given by $\left[\begin{array}{ccc}
x & 0 & 0   \\
 0 & y & 0 \\ 0 & z & x    
\end{array} \right]$.  Thus, the $\mathbb{R}^{3\times 3}$ transpose operation maps outside the algebra's representation.  However, there is an epimorphism to the two-dimensional subalgebra defined by elements $(x,y,0)$, and the latter's left regular representation is 
 $\left[\begin{array}{ccc}
x & 0  \\
 0 & y  
\end{array} \right]$.  The usual norm of the subalgebra is thus $xy$, and one of the special unital norms on the subalgebra is $(xy)^{\frac{1}{2}}$, and these differ only by an exponent - consistent with Lemma \ref{5/13/22.1}.   Note that the latter unital norm on the subalgebra provides the values for a special unital norm on the full algebra, corresponding to the choice of normalized uncurling metric in (\ref{11/17/22.2}) given by $\sigma = 0$, i.e., $\ell_0(x,y,z) = (xy)^{\frac{1}{2}}$, consistent with \mbox{Corollary \ref{11/17/22.1}}.  The fact that the transpose operation on $M_3(\mathbb{R})$ does not map the above left regular representation of $\left\{\left[\begin{array}{ccc}
x & z  \\
 0 & y  
\end{array} \right]:(x,y,z)\in \mathbb{R}^3\right\}$ to itself explains why there is no unital norm that is a power of the usual norm $x^2y$.
  
\subsection{Some higher dimensional algebras with unital norms that are sensitive to more element components than is  usual norm}

The usual norm on $ M_n(\mathbb{R})$ is sensitive to all components of an element of the algebra (because the determinant of an element of $M_n(\mathbb{R})$  is sensitive to the values of all entries of its matrix).    However, this is not true in general for associative algebras.  For example, for any algebra isomorphic to a triangular matrix algebra, the usual norm is only sensitive to the components on the main diagonal of the triangular matrix representation.  This is also true for the unital norm if all entries on the main diagonal of the triangular matrix representation are allowed to be distinct.  But if  at least two entries on the main diagonal always take equivalent values for the particular algebra,  then almost all of its unital norms family members will be sensitive to one or more components not on the main diagonal, in addition to all of the components on the main diagonal (these are matrix algebras for which an element that is not a diagonal matrix is never diagonalizable).  

It is easily seen from the presentation in Section 9 that the dual numbers are a two-dimensional example of this, and it
generalizes to higher dimensions as exemplified by Corollary \ref{5/17/22.1}.  We will now highlight an additional feature using the next three examples.  

\subsubsection{Another three-dimensional algebra $(\mathbb{T}_3)$}Consider the real algebra with elements  $(x,y,z)$ that multiply according to the matrices $ \left[
\begin{array}{ccc}
x & y & z   \\
0 & x & y \\
0 & 0 & x  
\end{array}
\right]$.    The constraints (\ref{7/21/23.1}) and (\ref{7/7/21.1}) lead to the normalized uncurling metrics, $ L = \left[
\begin{array}{ccc}
1 & \beta & \gamma   \\
\beta & \gamma & 0 \\
\gamma & 0 & 0  
\end{array}
\right]$, $\beta,\gamma\in\mathbb{R}$, which is just the $n=3$ case of Corollary \ref{7/1/22.1}.    Evaluating (\ref{11/7/20.2}), we now have a 2-parameter family of special unital norms, \begin{equation}\label{3/31/21.3} \ell_{\beta,\gamma}(x,y,z) = x\, e^{ \beta\frac{y}{x}  +\gamma \frac{z}{x} -\gamma \frac{ 1}{2}\left(\frac{y}{x}\right)^2 }.\end{equation}    The choice $\beta = \gamma = 0$ produces the special unital norm, $\ell_{0,0}(x,y,z)= x$, which is comparable to the usual norm, $x^3$.  
 The other special unital norms are distinguished by their dependence on $y$ and/or $z$, i.e., unlike the usual norm.  The non-units are associated with a different type of unital norm singularity (as $x$ approaches 0) than seen in examples considered so far other than $\mathbb{T}_n$.
 
 
 \subsubsection{A four-dimensional algebra} Consider the real unital algebra with elements 
$(x,y,z,w)$ which multiply according to the matrices 
$\left[
\begin{array}{ccccc}
x & y & w   \\
0 & x & z \\
0 & 0 & x  
\end{array}
\right]$.
Straightforward application of the constraints  (\ref{7/21/23.1}) and (\ref{7/7/21.1}) leads to the normalized uncurling metric family 
$\left[
\begin{array}{cccc}
1 & \beta & \gamma & \delta   \\
\beta & 0 & \frac{1}{2}\delta & 0 \\
\gamma & \frac{1}{2}\delta &  0 & 0   \\
\delta & 0 &  0 & 0
\end{array}
\right]$, and the associated special unital norm family,
\[ \ell_{\beta,\gamma,\delta}(x,y,z,w) = x\, e^{\beta\frac{y}{x} + \gamma\frac{z}{x} +\delta\frac{w}{x}- \frac{\delta}{2}\left(\frac{yz}{x^2}\right)}.\]


\subsubsection{A five-dimensional algebra}Consider the real unital algebra with elements 
$(x,y,z,v,w)$ which multiply according to the matrices
$\left[
\begin{array}{ccccc}
x & z & w   \\
0 & x & v \\
0 & 0 & y  
\end{array}
\right]$.
Straightforward application of the constraints (\ref{7/21/23.1}) and (\ref{7/7/21.1}) leads to the normalized uncurling metric family 
$\left[
\begin{array}{ccccc}
1+\sigma & 0 & \beta & 0  & 0   \\
0 & 1-\sigma & 0 & 0 & 0 \\
\beta & 0 &  0 & 0 & 0  \\
 0 & 0 &  0 & 0 & 0  \\
0 & 0 &  0 & 0 & 0
\end{array}
\right]$, and the associated special unital norm family 
\[ \ell_{\sigma,\beta}(x,y,z,v,w) =  (xy)^{\frac{1}{2}}\left(\frac{x}{y}\right)^{\frac{\sigma}{2}} e^{\frac{\beta}{2} \frac{z}{x}}.\]

In this example, almost all unital norms are influenced the element's $z$-component, while this sensitivity is not present with the usual norm (given by $x^4y$).  However, like the usual norm,  the unital norm remains insensitive to the $w$ and $v$ components.  Once again, there is a new unital norm singularity type as non-units are approached.  

Finally, the left regular representation of this algebra with elements $(x,y,z,v,w)$ is given by $ \left[
\begin{array}{ccccc}
x & 0 & 0 & 0 & 0   \\
0 & y & 0 & 0 & 0 \\
z & 0 &  x & 0 & 0  \\
0 & v &  0 & x & 0  \\
0 & w &  0 & z & x
\end{array}
\right].$ 
The transpose operation on $M_5(\mathbb{R})$ does not map this representation algebra into itself, due to its effect on the last three columns and last three rows - which explains why there is no unital norm that is a power of the usual norm $x^4y$.  However, there is an epimorphism to the two-dimensional subalgebra with elements $(x,y,0,0,0)$.  The subalgebra's left regular representation is $\left[\begin{array}{ccc}
x & 0  \\
 0 & y  
\end{array} \right]$, which leads to the special unital norm $(xy)^\frac{1}{2}$ according to \mbox{Lemma \ref{5/13/22.1}}.  and Theorem \ref{11/17/22.1}.

\vspace{.5cm}

Table III collects various examples presented in this section and earlier.  They illustrate the theorems of this paper in general. 
In particular, the factor in the expression for a special unital norm family that is devoid of a (Greek letter) parameter is the special unital norm family member that is closely allied with the usual algebra norm.   Perusal of the special unital norm family column also indicates algebras of the same dimension that cannot be isomorphic, and algebras of different dimensions (even if one is a subalgebra of another) that cannot be related by an epimorphism.

\hspace{-3.5cm} $ 
\begin{array}{cccccccc}
{\rm \underline{algebra}} & \mbox{\underline{usual norm}} & \mbox{\underline{special unital norm family}} & \mbox{\underline{anti-rotor ($M_{\mathfrak{u}}$) }} &  \bigskip \\ 
M_n(\mathbb{R}) & ({\rm det}(s))^n & ({\rm det}(s))^{\frac{1}{n}} & \alpha \,(n^2\times n^2)\mbox{\rm -``transpose\, matrix"}  \\  \\ 
\mathbb{C} & x^2+y^2 = r^2 & \sqrt{x^2+y^2}\,e^{\beta{\rm arctan}\left(\frac{y}{x}\right)} = re^{\beta\theta} & \left[\begin{array}{ccc}
 \alpha & \beta   \\ 
 \beta & -\alpha    
\end{array}\right]  =  \slashed{\mathbf{1}}_2 \, \mathbb{C}_\text{left} \\ \\ \slashed{\mathbb{C}}
 & x^2-y^2 & \begin{array}{cc}
 \sqrt{x^2-y^2}\, e^{\beta{\rm arctanh}\left(\frac{y}{x}\right)}   \\
 = (x^2-y^2)^{\frac{1}{2}}\left( \frac{x+y}{x-y}\right)^\frac{\beta}{2}    
\end{array} & \left[\begin{array}{ccc}
 \alpha & \beta   \\
 \beta & \alpha    
\end{array}\right]  = \slashed{\mathbb{C}}_\text{left} \\ \\
\mathbb{R}\times\mathbb{R} & xy & (xy)^{\frac{1}{2}}\left(\frac{x}{y}\right)^{\frac{\sigma}{2}} & \left[\begin{array}{ccc}
\sigma_1 & 0   \\
 0 & \sigma_2    
\end{array}\right] = [\mathbb{R}\times\mathbb{R}]_\text{left}  \\  \\
\vspace{-1mm} \mathbb{D}  & x^2 & x\, e^{\beta\frac{y}{x}} & \left[\begin{array}{ccc}
 \alpha & \beta   \\
 \beta & 0    
\end{array}\right] = \slashed{\mathbf{1}}_2\mathbb{D}_\text{left} \\ \\

\vspace{-1mm} \mathbb{H}  & (x^2+y^2+z^2+w^2)^2 & \sqrt{x^2+y^2+z^2+w^2} &  \,{\rm diag}\{\alpha,-\alpha,-\alpha,-\alpha\}   \\ \\

\mbox{\rm Cayley-Dixon seq., } \mbox{after }\mathbb{H} & (\mbox{\rm Euclidean norm})^2 & {\rm Euclidean\, Norm} &  \,{\rm diag}\{\alpha,-\alpha,-\alpha,\dots,-\alpha\} &  \\ \\
\mbox{Spin Factor Jordan} & (\mbox{\rm Minkowski norm})^2 & \mbox{\rm Minkowski norm} &  \,{\rm diag}\{\alpha,\alpha,\dots,\alpha\} & \\ \\ \prod_{i=1}^n \mathbb{R}
& \prod x_i & \left(\prod x_i^{\sigma_i}\right)^\frac{1}{n}, \sum\sigma_i=n & {\rm diag} \{\sigma_1,\dots,\sigma_n\} & \\ \\
\left\{ \left[\begin{array}{ccc}
 x & z   \\
 0 & y    
\end{array}\right]: (x,y,z)\in \mathbb{R}^3 \right\} & x^2y & (xy)^{\frac{1}{2}}\left(\frac{x}{y}\right)^{\frac{\sigma}{2}} & \left[\begin{array}{ccc}
\sigma_1  & 0 & 0   \\
 0 & \sigma_2 & 0 \\
 0 & 0 & 0    
\end{array}\right] & \\ \\
\left\{  \left[\begin{array}{ccc}
x  & y & z   \\
 0 & x & y \\
 0 & 0 & x    
\end{array}\right]: (x,y,z)\in \mathbb{R}^3   \right\}& x^3 & x\, e^{ \beta\frac{y}{x}  +\gamma \frac{z}{x} -\gamma \frac{ 1}{2}\left(\frac{y}{x}\right)^2 } &  \left[\begin{array}{ccc}
\alpha  & \beta & \gamma   \\
 \beta & \gamma & 0 \\
 \gamma & 0 & 0    
\end{array}\right]  \\ \\ \left\{ \left[
\begin{array}{ccccc}
x & y & w   \\
0 & x & z \\
0 & 0 & x  
\end{array}
\right] :  (x,y,z,w)\in\mathbb{R}^4 \right\} & x^4 & x\, e^{\beta\frac{y}{x} + \gamma\frac{z}{x} +\delta\frac{w}{x}- \frac{\delta}{2}\left(\frac{yz}{x^2}\right)}
& \left[
\begin{array}{cccc}
\alpha & \beta & \gamma & \delta   \\
\beta & 0 & \frac{1}{2}\delta & 0 \\
\gamma & \frac{1}{2}\delta &  0 & 0   \\
\delta & 0 &  0 & 0
\end{array}
\right]\\ \\
\left\{  \left[\begin{array}{ccc}
x  & z & w   \\
 0 & x & v \\
 0 & 0 & y    
\end{array}\right]:  (x,y,z,v,w)\in\mathbb{R}^5 \right\}  & x^4y & (xy)^{\frac{1}{2}}\left(\frac{x}{y}\right)^{\frac{\sigma}{2}}  e^{\frac{\beta}{2}  \frac{z}{x}}  & \left[\begin{array}{ccccc}
\sigma_1  & 0 & 0 & \beta & 0   \\
 0  & \sigma_2 & 0 & 0 & 0  \\
0   & 0 & 0 & 0 & 0 \\
 \beta  & 0 & 0 & 0 & 0  \\
 0  & 0 & 0 & 0 & 0    
\end{array}\right]  \\  \mathbb{T}_n & x_1^n & x_1 e^{P_{\gamma_2,\dots,\gamma_n}\left(\frac{x_2}{x_1},\dots,\frac{x_n}{x_1}\right)} & \slashed{\mathbf{1}}_n[\mathbb{T}_n]_\text{left} & \bigskip \end{array} 
$

\bigskip\bigskip

\enlargethispage{1cm}
\noindent {\bf Table III.} {\small Miscellaneous examples. $\prod_{i=1}^n \mathbb{R}$ is the direct product of rings.     $\mathbb{T}_n$ is the $n$-dimensional upper triangular Toeplitz matrix algebra.  $\slashed{\mathbf{1}}_n$ is the exchange matrix.  An expression with Greek letters is understood to imply the set whose members result from all realizations of those parameters. }

\newpage

\section{A conjecture}
As was remarked at the outset, we have resolved to study algebras by examining their implied nonlinear unary operations.  Based on results in Section 4 and Section 8,  we have initially proceeded from Definition \ref{7/21/23.1}.  
In fact, there is is some evidence that consideration of unary operations other than the multiplicative inverse might not add additional information regarding algebra structure.
 
  Specifically, alternative anti-rotors can (for example) be defined by replacing $s^{-1}$ with $s^j$, for $j\in\mathbb{Z}$, $j\ne 0,1$ in Definition \ref{7/21/23.1}, each referred to as ``type $j$".  For each of the algebras listed in the first column of Table III, these variant types of anti-rotors are easily computed by hand when  $|j|$ is ``small" (and $n$ is ``small" if applicable).  
Thereby one quickly discovers that, within each algebra, those anti-rotors of different type are {\it identical}.  

This result seems surprising.  That is, let $\hat{s}$ be the left regular representation of $s$, so that we can write $s^j = (\hat{s})^{j-1}[s]$, where $s$ ranges over the algebra elements in an appropriate neighborhood of $\mathbf{1}$. With respect to the basis defined by the columns of an uncurling metric $L$, the behavior described in the prior paragraph implies that the various vector fields $s\rightarrow s^j = (\hat{s})^{j-1}[s]$ are simultaneously irrotational, at least for small $|j|$ and $n$.  Now consider the vector field $F(s)[s]$ where $F(s)$ is an $(n\times n)$-matrix whose entries are smooth functions of components of $s$, and suppose $F(s)[s]$ is \mbox{irrotational}.  In general, this does {\it not} imply that $(F(s))^j[s]$ is irrotational for $j\ne 0,1$. For example, let $s=\left(\begin{array}{ccc}
 x   \\ 
y    
\end{array}\right)$, and $F= \left[\begin{array}{ccc}
x & 1   \\ 
 1 & y    
\end{array}\right]$.  Then $F(x,y)\left(\begin{array}{ccc}
 x   \\ 
y    
\end{array}\right)$ is irrotational.  But it is easily checked that $(F(x,y))^2\left(\begin{array}{ccc}
 x   \\ 
y    
\end{array}\right)$ is {\it not} irrotational.  So there is something special about the map defined by $\hat{s}$ that results in the unexpected behavior cited above. 
In practical terms, that behavior is of additional interest because it impacts the potential complexity of our general approach (i.e., the desirability of computing more than one anti-rotor type).  This important consideration is identified as a matter for future investigation. At this point we are simply left with,
\begin{conjecture} For a given algebra, the anti-rotors derived from any nonconstant nonlinear implied unary operation are identical. 
\end{conjecture}
\noindent The conjecture is a particularly emphatic assertion that nonlinear nonconstant unary operations each carry the same information regarding an algebra's structure, at least in the context of anti-rotors. 

\enlargethispage{1.5cm}
It is worthwhile observing that no matter how analysis of the Conjecture turns out, the result will be interesting.  \begin{enumerate}\item If the Conjecture is proved to be true, then the nonconstant nonlinear unary operations implied by the algebra product are simultaneously irrotational (i.e., an uncurling metric for one such operation uncurls all such operations) - which would seem to be a remarkable property of algebras. \item If the Conjecture is proved false, it becomes likely that there will be infinitely many classes of invariants derivable from this program, resulting from the infinitely many nonlinear unary operations associated with an algebra.  That is, the power of this approach would be accordingly much increased.  The alternative, that there are only finitely many nonlinear unary operations yielding distinct invariants, would simply be a variant of (1) above, and again represent a subtle feature of algebra structure.  \item If it is very difficult to settle the Conjecture (which would also mean that there is then a great deal of computer generated evidence in its favor), it would imply that the underlying property in question is a deep one. \end{enumerate}

\section{The case of nonunital algebras} Nonunital algebras can be addressed in a modification of Definition \ref{7/21/23.1} where the reference to an algebra's units is removed, and $s^{-1}$ is replaced by, e.g., $s^j$, $j>1$.   However, an invertible integral transform and associated dual group of unital norms will not always result.  

For nilpotent algebras, the anti-rotors of different types will be different - which is not the case for unital algebras if the Conjecture in Section 11 is correct.  For example, consider the real three-dimensional nonunital algebra composed of nilpotent elements $(x,y,z)$ that multiply according to the representation
  $ \left[
\begin{array}{ccccc}
0 & x & y & z  \\
0 & 0 & x & y \\
0 & 0 & 0 & x \\
0 & 0 & 0 & 0
\end{array}
\right] $.  We then have $(x,y,z)^2 = (0,x^2, 2xy)$.  The anti-rotor resulting from replacement of $s^{-1}$ with  $s^2$ in Definition \ref{7/21/23.1} is then $\left[
\begin{array}{ccccc}
\alpha & \beta & \gamma   \\
\beta & \gamma & 0 \\
\gamma & 0 & 0  
\end{array}
\right] $.  On the other hand, substitution of $s^{-1}$ with $s^3$  deals with $(x,y,z)^3 = (0,0, x^3)$, and the associated anti-rotor becomes $\left[
\begin{array}{ccccc}
\alpha & \beta & \gamma   \\
\beta & \delta & 0 \\
\gamma & 0 & 0  
\end{array}
\right] $.  Substitution of $s^{-1}$ by $s^j$, $j>3$, results in the additive group of all symmetric $(3\times 3)$-matrices since $s^j$ is zero for all algebra elements.  The anti-rotors associated with $s^2$, $s^3$, and $s^4$, must differ because the respective dimensions of the ranges of these operations differ - and this type of behavior will be typical with nilpotent algebras.  
 
The above case where $s^3$ replaces $s^{-1}$ in Definition \ref{7/21/23.1} is an example where the dimension of an anti-rotor exceeds the dimension of the algebra, and is relevant to comments made in Section 5.



\section{Issues for future development} 
It goes without saying that the homology functor has led to deep insights in the analysis of algebra structure.  In contrast, our ``anti-rotor program", which also addresses algebra structure and also has a functorial character, is embryonic  (though it has blossomed a bit from its Inverse Problems origins \cite{fgreensite:2022d}).  For example, the full extent of what can be gleaned from the Skewer, and especially the implied mapping of algebras to anti-rotors, has not been determined.  Collectively, \mbox{Corollary \ref{8/18/23.2}} and Corollary \ref{7/25/24.1} are surely only a partial listing of the invariants accessible with this approach.  More importantly, how (if at all possible) these might fit into a unified structure, as satisfying as a cohomology ring, say, is not at all clear.  So that vaguely stated challenge (a tall order, to say the least) is an overriding problem.  The following are more targeted and well defined questions.
 
\begin{enumerate}\item We have already identified the desirability of extending our program to \mbox{infinite} dimensional algebras in the context of characterizing the Hochschild cohomology ring for the case of algebras with infinite Hochschild cohomology dimension (Section 7).  If it proves possible to meaningfully provide such an extension, it would likely be of more general interest as well.
\item The essential feature of this approach is its differentiable aspect.  Thus, it is reasonable to suppose that extension or our program to complex algebras might be fruitful.  One simplifying aspect of that possibility can already be appreciated.  For real algebras, the multiplicative inversion operation is not uniquely forwarded for study as there are other nonlinear unary operations that are invertible (and thereby lead to possibly interesting integral transforms, i.e., duality, Section 3).  But in the context of complex algebras, the multiplicative inverse operation is uniquely forwarded as the only invertible nonlinear unary operation implied by the algebra axioms.
\item There are many potentially interesting combinatorial questions with a bearing on algebra structure that arise from the newly identified invariants in this paper (Section 5). 
\item What accessible structural information regarding an algebra does the image of the Skewer contain that is independent of the information obtainable from analysis of the kernel of the Skewer (Section 2)? 
\item Settle the Conjecture stated in Section 11.  It is likely that the more difficult this is, the greater will be its significance as regards algebra structure.
\item Is it possible to define an abelian anti-rotor category  (i.e., at least encompassing anti-rotors) on which a chain complex associated with an algebra arises, thus yielding a fourth family of isomorphism invariants (Section 7)?
\item Is the application to nonunital algebras worthwhile (Section 12)?
\end{enumerate}

\bigskip

\enlargethispage{1.5cm}

   \begin{thebibliography}{99}
\bibitem{eom:2024} Ringel CM, Encyclopedia of Mathematics (accessed June 13, 2024) Representation of an associative algebra. URL:http://encyclopediaofmath.org.
   
    \bibitem{drozd:1972} Drozd YA (1972) Representations of commutative algebras. Functional Analysis and its Applications, 6:286-288.
          
       \bibitem{belitskii:2005} Belitskii G, Lipyanski R, Sergeichuk VV (2005) Problems of classifying associative or Lie algebras and triples of symmetric or skew-symmetric matrices are wild.  Linear Algebra and its Applications, 407:249-262.
       
       \bibitem{hochschild:1945} Hochschild G (1945) On the cohomology groups of an associative algebra.  Annals of Mathematics, 46(1):58-67.
       
            \bibitem{gerstenhaber:1964} Gerstenhaber M (1964) On the deformation of rings and algebras, Annals of Mathematics, Second series, 79(1):59-103.
            
                     \bibitem{birkhoff:1932} Birkhoff, GD (1932), A Set of Postulates for Plane Geometry (Based on Scale and Protractors), Annals of Mathematics, 33: 329–345.
         
                       \bibitem{fgreensite:2022b} Greensite F (2022), A new proof of the Pythagorean Theorem and generalization of the usual algebra norm. arXiv:2209.14119

              \bibitem{bannwarth:2015} Bannwarth I, El Din M S (2015) Probabilistic algorithm for computing the dimension of real algebraic sets.  ISSAC'15 - 40th International Symposium on Symbolic and Algebraic Computation, July 2015, Bath UK, pp. 37-44, 10.1145/2755996.2756670, hal-01152751.
              
                 \bibitem{bourbaki:1989} Bourbaki N (1989) Algebra I, Springer, Berlin, p. 543.           
       
              \bibitem{kobayashi:2021} Kobayashi Y, Shirayanagi K, Tsukada M, Takahasi S (2021) A complete classification of three-dimensional algebras over $\mathbb{R}$ and $\mathbb{C}$. Asian-European Journal of Mathematics, 14(8) 2150131.                                        
   
              \bibitem{jacobson:1963} Jacobson N  (1963) Generic norm of an algebra. Osaka mathematics Journal, 15(1):25-50.
  
        \bibitem{schafer:1963}  Schafer R (1963) On forms of degree $n$ permitting composition. Journal of Mathematics and Mechanics, 12:777-792. 
         
      
       \bibitem{mccrimmon:2004} McCrimmon K (2004) A taste of Jordan algebras, Springer, New York.
       
                           \bibitem{fgreensite:2022d} Greensite F (2022), Origin of the uncurling metric formalism from an Inverse Problems regularization method. arXiv:2209.14137

              
      \end{thebibliography}

     \end{document}
     
     Thus, we have a set of points comprising a suitable unit ball around $\mathbf{1}$.   Each point has a norm $\ell_{\rm sp}(s)$ and a direction $\|\mathbf{1}\|^2 L_{\shortparallel}^{-1}[\nabla\ell_{\rm sp}(s^{-1})]$ as a point of the ``unit sphere" defined by $\{s\in\eta: \ell_{\rm sp}(s) = 1\}$.  By analogy with the definition of hyperbolic angle determined by the unit hyperbola in $\mathbb{R}^2$, we can then assign a signed angle to each point $s\in\eta$ as follows.  Consider the ordered triple of points $[(0,0),s, \mathbf{1}]$.  These define a planar subset of $\mathbb{R}^n$.    Orthogonally project the line segment with endpoints $\mathbf{1}$ and $s$ to the intersection of the ``unit sphere" with this plane.  In the plane defined by the three points of $[(0,0), s, \mathbf{1}]$, compute the area of the region bounded by the ray from $(0,0)$ to $s$, the projected curve on the portion of the ``unit sphere" lying in this plane, and the ray from $(0,0)$ to $\mathbf{1}$.  This area is assigned to be positive or negative according the right-hand-rule applied to the ordered triple of points $[(0,0),s, \mathbf{1}]$.  This signed area is the angle assigned to $s$.  The angle formed by $s_1$ and $s_2$ is then the difference of their individual assigned angles.

The point here is not just that we have exhibited a point set, and a norm and angle assigned to each point (which could be done in the above manner for normed linear spaces generally).  The idea is that we have a ``geometry" when a point can be expressed as the product of its ``norm'' and ``direction" with the latter being a point on the ``unit sphere". 

By analogy, comparing (\ref{5/22/24.5}), (\ref{5/22/24.6}) with (\ref{11/29/23.1}), (\ref{11/29/23.2}), we might say that distinct ``geometries" result from choices of $L_{\shortparallel}$ taken from different subspaces of $\mathfrak{u}$, since the latter endow an element with a norm and a direction that combine in a manner reminiscent of how a member of a quadratic space is expressed by a norm and direction.   More specifically, what we mean is that
